% concatenated from quant_oppenheim.tex
\documentclass[12pt]{amsart}
%\documentclass[11pt,a4paper]{amsart}
%\documentclass[11pt,leqno,a4paper]{amsart}
\usepackage[babel]{csquotes}
\usepackage[shortlabels]{enumitem}
\usepackage{amsmath,amsthm,amssymb,mathrsfs,amsfonts,verbatim,enumitem,color,leftidx,mathtools}
\usepackage{mathabx}
\usepackage{etoolbox} 
\usepackage{bbm}
\usepackage[all,tips]{xy}
\usepackage{graphicx,ifpdf}
\usepackage{stmaryrd}
\usepackage{amssymb}
\usepackage{dsfont}
% \usepackage{enumerate}
%\usepackage{bm}

\usepackage{geometry}
 \geometry{
 a4paper,
 total={160mm,247mm},
 left=25mm,
 top=25mm,
 }



\usepackage{color}
\definecolor{red}{rgb}{1,0,0}
\def\red{\color{red}}
\definecolor{gre}{rgb}{0,0.7,0}
\def\gre{\color{gre}}
\definecolor{blu}{rgb}{0,0,1}
\def\blu{\color{blu}}




\ifpdf
   \DeclareGraphicsRule{*}{mps}{*}{}
\fi
%%%%%%%%These package should be deleted before submitting
%\usepackage{hyperref}
%\usepackage[right,displaymath,mathlines]{lineno}
%\usepackage[notref,notcite]{showkeys} 
%\linenumbers
%%%%%%%%%%%%%%%%%%%%%%%%%%%%%%%%
\usepackage[colorlinks]{hyperref}
\hypersetup{
linkcolor=blue,          % color of internal links (change box color with linkbordercolor)
citecolor=green,        % color of links to bibliography
}

\usepackage{tikz}%pictures drawn in latex


\newtheorem{thm}{Theorem}[section]
\newtheorem{lem}[thm]{Lemma}
\newtheorem{cor}[thm]{Corollary}
\newtheorem{prop}[thm]{Proposition}
\newtheorem{prob}[thm]{Problem}

\theoremstyle{definition}
\newtheorem{defi}[thm]{Definition}
\newtheorem{ex}[thm]{Example}
\newtheorem{xca}[thm]{Exercise}

\theoremstyle{remark}
\newtheorem{remark}{Remark}[section]
\newtheorem{assume}{Assumption}

\numberwithin{equation}{section}
%\newtheorem{acknowledgments}{Acknowledgments}
% Absolute value notation
%\newcommand{\abs}[1]{\lvert#1\rvert}
%
%\makeatletter
%\newcommand{\rmnum}[1]{\romannumeral #1}
%\newcommand{\Rmnum}[1]{\expandafter\@slowromancap\romannumeral #1@}
%\makeatother

\definecolor{esperance}{rgb}{0.0,0.5,0.0}

\newcommand{\seon}[1]{\textsc{} {\color{magenta}\marginpar{} \textsf{#1}}}
\newcommand{\uri}[1]{\textsc{} {\color{blue}\marginpar{} \textsf{#1}}}
\newcommand{\nicolas}[1]{\textsc{} {\color{esperance}\marginpar{} \textsf{#1}}}

\newcommand{\ba}{\mathbf{a}}
\newcommand{\bb}{\mathbf{b}}
\newcommand{\bc}{\mathbf{c}}
\newcommand{\bd}{\mathbf{d}}
\newcommand{\be}{\mathbf{e}}
%\newcommand{\bf}{\mathbf{f}}
\newcommand{\bg}{\mathbf{g}}
\newcommand{\bh}{\mathbf{h}}
\newcommand{\bi}{\mathbf{i}}
\newcommand{\bj}{\mathbf{j}}
\newcommand{\bk}{\mathbf{k}}
\newcommand{\bl}{\mathbf{l}}
\newcommand{\bm}{\mathbf{m}}
\newcommand{\bn}{\mathbf{n}}
\newcommand{\bo}{\mathbf{o}}
\newcommand{\bp}{\mathbf{p}}
\newcommand{\bq}{\mathbf{q}}
%\newcommand{\br}{\mathbf{r}}
\newcommand{\bbk}{\mathbf{k}}

\newcommand{\bs}{\mathbf{s}}
\newcommand{\bt}{\mathbf{t}}
\newcommand{\bu}{\mathbf{u}}
\newcommand{\bv}{\mathbf{v}}
\newcommand{\bw}{\mathbf{w}}
\newcommand{\bx}{\mathbf{x}}
\newcommand{\by}{\mathbf{y}}
\newcommand{\bz}{\mathbf{z}}

\newcommand{\balpha}{\boldsymbol{\alpha}}
\newcommand{\bbeta}{\boldsymbol{\beta}}
\newcommand{\bxi}{\boldsymbol{\xi}}
\newcommand{\bupsilon}{\boldsymbol{\upsilon}}


\newcommand{\lefttext}[1]{\makebox[0pt][l]{#1}}
\newcommand{\fP}{\mathfrak{P}}
\newcommand{\fB}{\mathfrak{B}}

\newcommand{\dd}{\mathrm{d}}

\newcommand{\F}{\mathbb{F}}
\newcommand{\R}{\mathbb{R}}
\newcommand{\N}{\mathbb{N}}
\newcommand{\Z}{\mathbb{Z}}
\newcommand{\Q}{\mathbb{Q}}
\newcommand{\C}{\mathbb{C}}
\newcommand{\EE}{\mathbf{e}}
\newcommand{\sm}{\smallsetminus}
\newcommand{\til}{\widetilde}
\newcommand{\ve}{\varepsilon}
\newcommand{\ep}{\varepsilon}
\newcommand{\s}{\sigma}
\newcommand{\de}{\delta}

\newcommand{\G}{\Gamma}

\newcommand{\e}{\mathrm{e}}

\newcommand{\Yb}{Y_{\on{bad}}}
\newcommand{\Ybe}{F_{S^+(\cO)}}
\newcommand{\Ybr}{E^\cO_{[r,\infty)}}
\newcommand{\Ybm}{\E{r_m}}
\newcommand{\Ybmprime}{E^{\cO'}_{[r_m,\infty)}}
%{Y_{\on{bad}}^{\cO,r_m}}
\newcommand{\Ybee}{Y_{\on{bad}}^{{\eps^\prime}}}
\newcommand{\Ybem}{Y_{\on{bad}}^{{\eps^\prime,r_m}}}
\newcommand{\Ha}{\mathcal{H}}
%\newcommand{\wt}{\widetilde}}


%\newcommand{\diag}{\mathrm{diag}}
\newcommand{\spa}{\mathrm{span}}
\newcommand{\Ad}{\operatorname{Ad}}
\newcommand{\supp}{\mathrm{supp\,}}
\newcommand{\rk}{\mathrm{rk\,}}
\newcommand{\A}{\cA}


%%%%%Redefine an existing command
%\renewcommand{\Latex}{\LaTex \xspace}
%%%%%%%%%%%%%%%%%%




% Blank box placeholder for figures (to avoid requiring any
% particular graphics capabilities for printing this document).
\newcommand{\blankbox}[2]{%
  \parbox{\columnwidth}{\centering
% Set fboxsep to 0 so that the actual size of the box will match the
% given measurements more closely.
    \setlength{\fboxsep}{0pt}%
    \fbox{\raisebox{0pt}[#2]{\hspace{#1}}}%
  }%
}



%              Latin Letters
%------------------------------------------------------------------------------------------------------------
\newcommand{\al}{\alpha}
%\newcommand{\be}{\beta}
\newcommand{\ga}{\gamma}
\newcommand{\Ga}{\Gamma}
\newcommand{\del}{\delta}
\newcommand{\Del}{\Delta}
\newcommand{\lam}{\lambda}
\newcommand{\Lam}{\Lambda}
\newcommand{\eps}{\epsilon}
\newcommand{\veps}{\varepsilon}
\newcommand{\ka}{\kappa}
\newcommand{\sig}{\sigma}
\newcommand{\Sig}{\Sigma}
\newcommand{\om}{\omega}
\newcommand{\Om}{\Omega}
\newcommand{\vphi}{\varphi}
%------------------------------------------------------------------------------------------------------------
%-----Capital letters--------------
\newcommand{\cA}{\mathcal{A}}
\newcommand{\cB}{\mathcal{B}}
\newcommand{\cC}{\mathcal{C}}
\newcommand{\cD}{\mathcal{D}}
\newcommand{\cE}{\mathcal{E}}
\newcommand{\cF}{\mathcal{F}}
\newcommand{\cG}{\mathcal{G}}
\newcommand{\cH}{\mathcal{H}}
\newcommand{\cI}{\mathcal{I}}
\newcommand{\cJ}{\mathcal{J}}
\newcommand{\cK}{\mathcal{K}}
\newcommand{\cL}{\mathcal{L}}
\newcommand{\cM}{\mathcal{M}}
\newcommand{\cN}{\mathcal{N}}
\newcommand{\cO}{\mathcal{O}}
\newcommand{\cP}{\mathcal{P}}
\newcommand{\cQ}{\mathcal{Q}}
\newcommand{\cR}{\mathcal{R}}
\newcommand{\cS}{\mathcal{S}}
\newcommand{\cT}{\mathcal{T}}
\newcommand{\cU}{\mathcal{U}}
\newcommand{\cV}{\mathcal{V}}
\newcommand{\cW}{\mathcal{W}}
\newcommand{\cX}{\mathcal{X}}
\newcommand{\cY}{\mathcal{Y}}
\newcommand{\cZ}{\mathcal{Z}}



\newcommand{\bA}{\mathbb{A}}
\newcommand{\bB}{\mathbb{B}}
\newcommand{\bC}{\mathbb{C}}
\newcommand{\bD}{\mathbb{D}}
\newcommand{\bG}{\mathbb{G}}
\newcommand{\bL}{\mathbb{L}}
\newcommand{\bM}{\mathbb{M}}
\newcommand{\bP}{\mathbb{P}}
\newcommand{\bR}{\mathbb{R}}
\newcommand{\bZ}{\mathbb{Z}}
\newcommand{\bQ}{\mathbb{Q}}
\newcommand{\bF}{\mathbb{F}}
\newcommand{\bK}{\mathbb{K}}
\newcommand{\bN}{\mathbb{N}}
\newcommand{\bH}{\mathbb{H}}
\newcommand{\bT}{\mathbb{T}}
\newcommand{\bS}{\mathbb{S}}




\newcommand{\goa}{\mathfrak{a}}
\newcommand{\gob}{\mathfrak{b}}
\newcommand{\goc}{\mathfrak{c}}
\newcommand{\gog}{\mathfrak{g}}
\newcommand{\goh}{\mathfrak{h}}
\newcommand{\gol}{\mathfrak{l}}
\newcommand{\gon}{\mathfrak{n}}
\newcommand{\gop}{\mathfrak{p}}
\newcommand{\gos}{\mathfrak{s}}
\newcommand{\got}{\mathfrak{t}}
\newcommand{\gou}{\mathfrak{u}}



%%%%%%%Groups
\newcommand{\SL}{\operatorname{SL}}
\newcommand{\SO}{\operatorname{SO}}
\newcommand{\PO}{\operatorname{PO}}
\newcommand{\PGL}{\operatorname{PGL}}
\newcommand{\ASL}{\operatorname{ASL}}
\newcommand{\GL}{\operatorname{GL}}
\newcommand{\PSL}{\operatorname{PSL}}
\newcommand{\Id}{\operatorname{Id}}
\newcommand{\Sp}{\operatorname{Sp}}

\newcommand{\vol}{\operatorname{vol}}
\newcommand{\dist}{\operatorname{dist}}
\newcommand{\ord}{\operatorname{ord}}


\newcommand*{\transp}[2][-1mu]{\ensuremath{\mskip1mu\prescript{\smash{\mathrm t\mkern#1}}{}{\mathstrut#2}}}
%------------------------------------------------------------------------------------------------------------
%              Special commans
%------------------------------------------------------------------------------------------------------------
%\newcommand{\def}{\overset{\on{def}}{=}}
\newcommand{\comp}{\textrm{{\tiny $\circ$}}}
\newcommand{\lie}{\operatorname{Lie}}
\newcommand\bignorm[1]{\left\|#1\right\|}
\newcommand\norm[1]{\|#1\|}
\newcommand\wt[1]{\widetilde{#1}}
\newcommand\wh[1]{\widehat{#1}}
\newcommand\set[1]{\left\{#1\right\}}
\newcommand\pa[1]{\left(#1\right)}
\newcommand\idist[1]{\langle#1\rangle}
\newcommand\bigav[1]{\left|#1\right|}
\newcommand\av[1]{|#1|}
\newcommand\on[1]{\operatorname{#1}}
\newcommand\diag[1]{\operatorname{diag}\left(#1\right)}
\newcommand\mb[1]{\mathbf{#1}}
\newcommand\E[1]{E^\cO_{[#1,\infty)}}


\newcommand\tb[1]{\textbf{#1}}
\newcommand\mat[1]{\pa{\begin{matrix}#1\end{matrix}}}
\newcommand\br[1]{\left[#1\right]}
\newcommand\smallmat[1]{\pa{\begin{smallmatrix}#1\end{smallmatrix}}}
\newcommand\crly[1]{\mathscr{#1}}
\newcommand\tran[1]{{}^t#1}
\newcommand{\re}{\operatorname{Re}}
\newcommand{\im}{\operatorname{Im}}
\newcommand{\Mat}{\operatorname{Mat}}
\newcommand{\pr}{\operatorname{pr}}
\newcommand{\wstar}{\overset{\on{w}^*}{\lra}}
%\newcommand{\spa}{\on{span}}
\newcommand{\Supp}{\on{Supp}}
\newcommand{\dsim}{\Del_{\on{sym}}}
\newcommand{\defn}{\overset{\on{def}}{=}}
%------------------------------------------------------------------------------------------------------------
%------------------------------------------------------------------------------------------------------------
%             Arrows
%------------------------------------------------------------------------------------------------------------
%------------------------------------------------------------------------------------------------------------

\newcommand{\lla}{\longleftarrow}
\newcommand{\la}{\leftarrow}
\newcommand{\lra}{\longrightarrow}
\newcommand{\ra}{\rightarrow}
\newcommand{\ora}[1]{\overset{#1}{\rightarrow}}
\newcommand{\olra}[1]{\overset{#1}{\longrightarrow}}
\newcommand{\olla}[1]{\overset{#1}{\longleftarrow}}
\newcommand{\ulra}[1]{\underset{#1}{\longrightarrow}}
\newcommand{\ulla}[1]{\underset{#1}{\longleftarrow}}
\newcommand{\hra}{\hookrightarrow}
\newcommand{\lera}{\leftrightarrow}
\newcommand{\onto}{\xymatrix{\ar@{>>}[r]&}}
% Im not sure what is the difference between the above \olra, \ora and the following
\newcommand{\slra}[1]{\stackrel{#1}{\lra}}
\newcommand{\sra}[1]{\stackrel{#1}{\ra}}

\newcommand{\dima}{\underline{\dim}_a}
\newcommand{\dimm}{\underline{\dim}_M}

%------------------------------------------------------------------------------------------------------------
%------------------------------------------------------------------------------------------------------------
%        Makros for equations enumeration etc.
\newcommand{\eq}[1]
{
\begin{equation*}
{#1}
\end{equation*}
}
%%%%%%%%%%%%
\newcommand{\eqlabel}[2]
{
\begin{equation}
{#2}\label{#1}
\end{equation}
}
%%%%%%%%%%%%%%
\newcommand{\eqalign}[1]
{
\begin{align*}
{#1}
\end{align*}
}
%%%%%%%%%%%%%%
\newcommand{\eqalignlabel}[2]
{
\begin{align}
{#2}\label{#1}
\end{align}
}
%%%%%%%%%%%%%%
%-----------------------------------------------    Notes of coauthors     
\newcommand{\usnote}[1]{\marginpar{\color{blue}\tiny [US] #1}}
\newcommand{\nsnote}[1]{\marginpar{\color{esperance}\tiny [NS] #1}}
\newcommand{\usadd}[1]{{\color{red}{\tiny [US]} #1}}
%%%
%%% it is possible to change the font from tiny to even smaller. here are two possibilities
% \font\sn = cmssi8 scaled \magstep0
% \font\ssn = cmssi5 scaled \magstep0
% To use them write {\sn ...text... } for example.
%-----------------------------------------------------
\newcommand{\seonnote}[1]{\marginpar{\color{magenta}\tiny [SS] #1}}


\makeatletter
\newcommand*{\rom}[1]{\expandafter\@slowromancap\romannumeral #1@}
\makeatother

\begin{document}

\title[Values of ternary quadratic forms at integers]{Values of ternary quadratic forms at integers and the Berry-Tabor conjecture for 3-tori}
\author{Wooyeon Kim \and Jens Marklof \and Matthew Welsh}

\address{Wooyeon Kim, June E Huh Center for Mathematical Challenges, Korea Institute for Advanced Study, Seoul, Republic of Korea \newline \rule[0ex]{0ex}{0ex} \hspace{8pt}{\tt wooyeonkim@kias.re.kr}}
\address{Jens Marklof, School of Mathematics, University of Bristol, Bristol BS8 1UG, U.K.\newline \rule[0ex]{0ex}{0ex} \hspace{8pt}{\tt j.marklof@bristol.ac.uk}}
\address{Matthew Welsh, Department of Mathematics, Michigan State University, 619 Red Cedar Road, C212 Wells Hall, East Lansing, MI 48864 U.S.A. \newline \rule[0ex]{0ex}{0ex} \hspace{8pt}{\tt welshmat@msu.edu}}



\date{6 January 2026}
\thanks{JM's research was supported by EPSRC grant EP/W007010/1. WK was supported by Korea Institute for Advanced Studies (HP101301).}

\begin{abstract}
Berry and Tabor conjectured in 1977 that spectra of generic integrable quantum systems have the same local statistics as a Poisson point process. We verify their conjecture in the case of the two-point spectral density for a quantum particle in a three-dimensional box, subject to a Diophantine condition on the domain's proportions. A permissible choice of width, height and depth is for example $1,2^{1/3},2^{-1/3}$. 
This extends previous work of Eskin, Margulis and Mozes (Annals of Math., 2005) in dimension two, where the problem reduces to the quantitative Oppenheim conjecture for quadratic forms of signature $(2,2)$. 
The difficulty in three and higher dimensions is that we need to consider the distribution of indefinite forms in shrinking rather than fixed intervals, which we are able to resolve for special diagonal forms of signature $(3,3)$ in various scalings, including a rate of convergence. A key step of our approach is to represent the relevant counting problem as an average of a theta function on $\SL(2,\Z)^3\backslash\SL(2,\R)^3$ over an expanding family of one-parameter unipotent orbits. The asymptotic behaviour of these unipotent averages follows from Ratner's measure classification theorem and subtle escape of mass estimates. 
\end{abstract}

\maketitle


%\tableofcontents

\section{Introduction}
\label{sec:intro}

\subsection{The Berry-Tabor conjecture}

A principle motivation for this work is the Berry-Tabor conjecture \cite{BerryTabor1977}, which asserts that the statistical correlations in the quantum spectrum of a generic, classically integrable system are that of a Poisson point process, i.e., are maximally uncorrelated. In contrast, the statistics of quantum systems with ``chaotic'' classical limit are expected -- according to Bohigas, Giannoni and Schmit \cite{BGS1984} -- to be governed by one of the Gaussian random matrix ensembles, where the choice of ensemble depends on the time-reversal symmetries of the system. Despite overwhelming numerical evidence and major theoretical advances, there is currently no complete proof of either phenomenon for a single example. 

We will here investigate the validity of the Berry-Tabor conjecture for a specific model, a particle confined to a $d$-dimensional box with edge lengths $\ell_1,\ldots,\ell_d$. The particle's classical motion is along straight lines with elastic reflection at the boundary. This motion lifts to the Kronecker flow on a $d$-torus obtained by reflection on all faces of the box, and is therefore completely integrable. The standard quantisation of this dynamical system is given by the Hamiltonian $H=-\Delta_D$, where $\Delta_D$ is the Dirichlet Laplacian on the cuboid $[0,\ell_1]\times\cdots\times[0,\ell_d]\subset\R^d$. 
The spectrum 
$0<\lambda_1\leq\lambda_2\leq \cdots \to\infty$ of the quantum Hamiltonian $H=-\Delta_D$ is given by the values of the positive definite quadratic form
\begin{equation}\label{Q0}
Q_0(m) = \sum_{j=1}^d  x_j m_j^2 , \qquad x_j = \frac{\pi^2}{\ell_j^2}, 
\end{equation}
where $m=(m_1,\ldots,m_d)$ runs over the integer vectors in $\Z^d$ with positive coefficients. 

The spectrum of the Laplacian for the above $d$-torus (comprising $2^d$ copies of the box) is represented by the same quadratic form, but $m$ now runs over all of $\Z^d$. Our results therefore apply directly to this case. In fact, the reflection symmetry of the rectangular torus under the dihedral group $D_{2h}$ allows a decomposition of its spectrum according to the one-dimensional unitary representations of $D_{2h}$. These are in turn realised by the Laplacian with Dirichlet or Neumann boundary conditions on the various faces of the original box.  

Weyl's law for the spectrum of $H=-\Delta_D$ tells us that
\begin{equation}\label{WeylsLaw}
\#\{ j : \lambda_j < T \} \sim \frac{\ell_1\cdots\ell_d}{(2\pi)^d} \,\omega_d \, T^{d/2} \qquad (T\to\infty),
\end{equation}
where $\omega_d$ denotes the volume of the unit ball in $\R^d$. In dimension $d=2$ the law is linear in $T$ and hence the average spacing between consecutive elements is asymptotically constant. This is no longer the case in dimension three and higher. Thus, in order to measure correlations on the correct local scale, we unfold the spectrum by setting
\begin{equation}\label{unfold}
\xi_j = \frac{\ell_1\cdots\ell_d}{(2\pi)^d} \,\omega_d \, \lambda_j^{d/2} .
\end{equation}
This simple rescaling produces a sequence whose average spacing is now asymptotic to constant one.

For $T>0$, the two-point (or pair correlation) measure $\cR_T$ of the unfolded, truncated sequence $(\xi_j \mid \xi_j<T)_j$ is defined by
\begin{equation}
\cR_T(A) = \#\{ (i,j) \mid i\neq j, \; \xi_i<T,\; \xi_j \in A+\xi_i\} ,
\end{equation}
with $A\subset\R$ a bounded Borel set. We are interested in the asymptotic behaviour for large $\xi$, and say that the two-point measure of the full sequence $(\xi_j)_j$ is Poissonian if
\begin{equation}
\lim_{\xi\to\infty} \frac{ \cR_T(A) }{T}  = |A| 
\end{equation}
for any Borel set $A\subset\R$ with $|\partial A|=0$. Here $|A|$ denotes the Lebesgue measure of $A$.
The reason for this terminology is that the above holds almost surely if $(\xi_j)_j$ is a realisation of a Poisson point process on $\R_{>0}$ of intensity one. 

We say that a real number $x$ is $\kappa$-Diophantine for some $\kappa \geq 1$ if there exists a finite constant $C>0$ such that $|x - \tfrac{r}{q}| \geq C q^{-\kappa -1}$ for all integers $r\in\Z$, $q\in\Z_{\neq 0}$. For example, all quadratic surds are $1$-Diophantine, and all irrational algebraic numbers are $\kappa$-Diophantine for every $\kappa>1$ (Roth's theorem). Furthermore, for any given $\kappa>1$, the set of $\kappa$-Diophantine numbers has full Lebesgue measure. 

The following theorem proves the Berry-Tabor conjecture for the second-order statistics of the spectrum of the Hamiltionan $H=-\Delta_D$. Equivalently, it states that the two-point correlations of the values of the degree $d$ homogeneous form $Q_0(m)^{d/2}$ at integers $m\in\Z_{\geq 1}^d$ are Possonian. 

\begin{thm}
  \label{thm:main}
  Let $(\xi_j)_j$ be the deterministic sequence given by the unfolded spectrum \eqref{unfold} of the Dirichlet Laplacian for a box in $\R^3$ with edge lengths $\ell_1,\ell_2,\ell_3$.
  If the ratios $\ell_i^2/\ell_j^2$ are $\kappa$-Diophantine for all $i \neq j$ with $\kappa$ sufficiently close to 1, then the pair correlation measure of $(\xi_j)_j$ is Poissonian.
 \end{thm}

 By Roth's theorem, the conditions of the theorem are satisfied, e.g., for $\ell_1= 1$ and algebraic numbers $\ell_2,\ell_3$ so that $\ell_2^2,\ell_3^2,\ell_2^2/\ell_3^2\notin\Q$. This includes the example given in the abstract. The analogous statement was proved by Eskin, Margulis and Mozes \cite{EskinMargulisMozes2005} for positive definite binary quadratic forms $Q_0$, under similar Diophantine conditions, where (as we will discuss further below) this is a special -- and particularly difficult -- case of the quantitative Oppenheim conjecture.
 The necessary unfolding for $d \geq 3$ requires a different approach from that taken in \cite{EskinMargulisMozes2005} or in the more recent effective version by Lindenstrauss, Mohammadi and Wang \cite{LMW2023b}.
 We will here use a similar strategy as in \cite{Marklof2002,Marklof2003} and translate the problem to convergence of unipotent averages over Siegel theta functions associated to $Q_0$. 
Theta functions were also used in connection with quantitative versions of the Oppenheim conjecture, as well as the Davenport-Lewis conjecture, for quadratic forms with five or more variables in the work of G\"otze \cite{Goetze2004} and Buterus, G\"otze, Hille and Margulis \cite{BGHM2022}. 

It is well known that the Diophantine conditions in Theorem \ref{thm:main} are satisfied for $\ell_1=1$ and almost every (in the sense of Lebesgue measure) choice of $\ell_2,\ell_3$. Our theorem therefore strengthens, in dimension $d=3$, recent results by Blomer and Li \cite{BlomerLi2023} proved for almost all rectangular tori.
Previous proofs of the Poisson nature of almost all flat tori are due to \cite{Sarnak1997} in dimension 2, and VanderKam \cite{VanderKam1999a,VanderKam1999b,VanderKam2000} in dimension $d\geq 4$. VanderKam \cite{VanderKam2000} also proves convergence to the Poissonian $k$-point correlation measure, provided $d\geq 2k$. 


Triple and higher correlation densities are much less understood in the deterministic setting, i.e., under explicit conditions on the edge lengths $\ell_i$. The best current results for binary forms are due to Aistleitner, Blomer and Radziwi\l\l\ \cite{ABR2024}. The problem of understanding the distribution of gaps $\xi_{j+1}-\xi_j$ appears to be equally difficult; cf.~Bourgain, Blomer, Rudnick and Radziwi\l\l\ \cite{BBRR2017}.

The Diophantine conditions in Theorem \ref{thm:main} are likely not sharp. Note, however, the statement evidently fails for all rational length ratios, when the two-point measure diverges. Following Sarnak \cite{Sarnak1997}, see also \cite{VanderKam1999a,VanderKam1999b,VanderKam2000,Marklof2002,Marklof2002b,Marklof2003}, one can in fact show that the divergence persists for a set of second Baire category of edge lengths, i.e., topologically generic boxes.


\subsection{Uniform variant of the quantitative Oppenheim conjecture}

Instead of unfolding the sequence as in \eqref{unfold}, we can also count values in shrinking intervals to compensate for the increasing mean density. This will lead to a uniform variant of the quantitative Oppenheim conjecture of Eskin, Margulis and Mozes \cite{EskinMargulisMozes1998,EskinMargulisMozes2005}, which we will prove here for the special quadratic forms of signature $(3,3)$ given by
\begin{equation}
Q(m) = Q_0(m_1)-Q_0(m_2), \qquad m=(m_1,m_2),
\end{equation}
where $Q_0$ is the positive definite ternary diagonal form in \eqref{Q0}.

Denote by $\cS(\R^n)$ the Schwartz class of functions $\R^n\to\C$. For $M,L>0$, $F\in\cS(\R^d)$, $\psi\in\cS(\R)$ we define the counting measure $\cN_{M,L}$ by
\begin{equation}
 \label{eq:NLMdef}
  \cN_{M,L}(F,\psi) 
  =  \sum_{m\in \Z^{2d}} F(M^{-1} m) \psi ( L Q(m)) ,
\end{equation}
and extend this to Borel measurable functions as usual. Thus, if $F$ and $\psi$ are characteristic functions of a unit ball and interval $[a,b]$, respectively, then $\cN_{M,L}$ describes the distribution of values of $Q(m)$ in the interval $[a/L,b/L]$ as $m$ runs over all lattice points $\Z^d$ in the ball of radius $M$. In the classical setting of the quantitative Oppenheim conjecture \cite{EskinMargulisMozes1998,EskinMargulisMozes2005}, $L$ is fixed. A key observation is that in the signature $(2,2)$ case Diophantine conditions on the coefficients are necessary, whereas for signature $(3,1)$ and indefinite forms in more than four variables, only irrationality of the form is required.
However, this changes
if we count in shrinking intervals, i.e., $L\to\infty$ as $M\to\infty$ at suitable rates.


In view of \eqref{WeylsLaw}, the mean spacing between the values of $Q_0(m)$ is about $M^{-(d-2)}$, and hence we expect $\cN_{M,L}(F,\psi)$ to be of size $\asymp M^d\times M^{d-2}/L=M^{2d-2}/L$. The scaling $L=M^{d-2}$ is the critical scaling which captures the local correlations relevant to Theorem \ref{thm:main}, and we will see below that the following statement in fact implies the former as a corollary. 

Let $\e(z)=\exp(2\pi\mathrm{i} z)$, and denote by $D_{2h}$ the symmetry group of a generic rectangular cuboid in $\R^d$, centred at the origin, with $\ord D_{2h}=2^d$. 

\begin{thm}
  \label{thm:main2}
  Fix $d=3$.
  If the ratios $x_i/x_j$ are $\kappa$-Diophantine for all $i \neq j$ with $\kappa$ sufficiently close to 1, then for $\psi\in\cS(\R)$, $F\in\cS(\R^{2d})$, $L,M\geq 1$,
\begin{multline}\label{eq:main2}
\cN_{M,L}(F,\psi)   
= \frac{M^{2d-2}}{L}  \int_{\R} \psi(x) \dd x \;\int_{\R} \bigg( \int_{\R^{2d}} F(v) \e( \tfrac{1}{2} \xi Q(v)) \dd v \bigg) \dd \xi \\
+ M^d \, \psi(0) \sum_{\sigma\in D_{2h}} \int_{\R^d} F((u,\sigma(u))) \dd u  + O_{F,\psi}\bigg(\frac{M^{2d-2-\delta}}{L}\bigg)+o_{F,\psi}(M^d).
\end{multline}
for any $\delta < \tfrac{1}{\kappa + 1}$. 
\end{thm}

We remark that for the special choice $F(u_1,u_2)=\varphi_1(Q_0(u_1)) \varphi_2(Q_0(u_2))$ with $\varphi_i\in\cS(\R_{\geq 0})$, \eqref{eq:main2} becomes
\begin{multline}\label{eq:main222}
\cN_{M,L}(F,\psi)   
= \frac{M^{2d-2}}{L} \, \frac{d^2}{2}\, \frac{\omega_d^2}{x_1\cdots x_d} \int_{\R} \psi(x) \dd x \; \int_0^\infty \varphi_1(\lambda) \varphi_2(\lambda) \lambda^{d-2} d\lambda  \\
+ 2^dM^d \,\frac{d}{2}\, \frac{\omega_d}{(x_1\cdots x_d)^{1/2}}  \psi(0) \int_{\R^d} \varphi_1(\lambda)\varphi_2(\lambda) \lambda^{\frac{d}{2}-1} \dd \lambda  + O_{F,\psi}\bigg(\frac{M^{2d-2-\delta}}{L}\bigg)+o_{F,\psi}(M^d).
\end{multline}
If $L=M^{d-2}$, Theorem \ref{thm:main} now follows from \eqref{eq:main222} by the approximation argument of \cite[Section 2]{Marklof2002}. (Note that we need to take into account the dihedral symmetry, which explains the additional factor of $2^d$.)



% We conjecture that Theorem \ref{thm:main2} holds also for $d\geq 4$ under suitable Diophantine conditions.
As remarked earlier, Eskin, Margulis and Mozes established \eqref{eq:main2} (without a rate) for $L$ fixed and $d\geq 3$ in \cite{EskinMargulisMozes1998} only assuming at least one of the ratios $x_i/x_j$ is irrational, and for $d=2$ \cite{EskinMargulisMozes2005} allowing $\kappa$-Diophantine ratios for any $\kappa\geq 1$.
Also for $d = 2$ and $\kappa$-Diophantine ratios, Lindenstrauss, Mohammadi and Wang have proved a version of (\ref{eq:main2}) with an effective bound for the remainder, allowing for $L$ to grow as a small power of $M$.

The approach taken in \cite{EskinMargulisMozes2005} requires $L$ fixed and the bound for remainder in \cite{LMW2023b} depends quite sensitively on $L$, so taking $L$ as large as $M^{d-2}$ for $d \geq 3$ requires a different strategy.
Here we use Fourier inversion on $\psi$ to transform $\cN_{M,L}(F,\psi)$ into a unipotent average of a theta function.
The advantage is that $L$ appears only in the length of the unipotent average and not implicitly in the test functions.
This feature is used in a similar way in \cite{Marklof2002,Marklof2003} and as well as by Buterus, G\"otze, Hille and Margulis \cite{BGHM2022}. 

Buterus, G\"otze, Hille and Margulis \cite{BGHM2022} obtain convergence rates of the form \eqref{eq:main2} for general non-singular quadratic forms (both definite and indefinite) with more than five variables for expanding ($L\to 0$) and shrinking intervals ($L\to\infty$), as $M\to\infty$ under Diophantine conditions similar to those in \cite{EskinMargulisMozes2005}.
While our work provides a stronger rate of convergence in the special case we consider,
%under the strongest Diophantine condition, they obtain $ L^{\frac{1}{18}} M^{\frac{34}{9}}$, smaller than $L^{-1} M^4$ if $L < M^{\frac{4}{19}}$, while we get $L^{-1} M^{\frac{5}{2}}$)***
the key breakthrough is that we can deal with the scaling regime $M,L\to\infty$ with $L\geq M^{d-2-\delta}$, which is when the second term on the right hand side of \eqref{eq:main2} dominates the error term. This includes the critical scaling $L= M^{d-2}$ required for the proof of Theorem \ref{thm:main}.
A further advantage of our approach, particularly relevant to the calculation of the second term, is the deployment of theta sums with general cut-off functions $F$, via the Shale-Weil representation, not only classical theta series with Gaussian weights as in \cite{BGHM2022}.

As mentioned above, a rate of convergence for signature $(2,2)$, in the critical scaling ($L$ fixed), was recently obtained by Lindenstrauss, Mohammadi and Wang \cite{LMW2023b} as an application of their ground-breaking work on effective Ratner equidistribution \cite{LMW2022,LMW2023a}. 
Str\"ombergsson and Vishe \cite{AV2020} obtained rates of convergence for the quantitative Oppenheim conjecture for special inhomogeneous forms studied in \cite{Marklof2002,Marklof2003}. The quantitative Oppenheim conjecture for general inhomogeneous forms was proved by Mohammadi and Margulis \cite{MargulisMohammadi2011} without rate of convergence. More is known if one restricts to error rates valid for almost every quadratic form of a given signature; see the papers by Bourgain \cite{Bourgain2016}, Ghosh and Kelmer \cite{GhoshKelmer2017,GhoshKelmer2018}, Athreya and Margulis \cite{AthMar2018}, and Ghosh, Kelmer and Yu \cite{GKY2022,GKY2023}.


\subsection{Escape of mass and Margulis functions}

Our strategy for the proof of Theorem \ref{thm:main2} is to relate $\cN_{M,L}(F,\psi)$ to homogeneous dynamics on $\Gamma^d \backslash G^d$, $\Gamma = \SL(2, \Z)$, $G = \SL(2, \R)$, via a theta function $\widetilde\Theta_F$ defined on the semi-direct product of a Heisenberg group with $G^d$ using the Shale-Weil representation. The connection between the counting function $\cN_{M,L}(F,\psi)$ and theta function $\widetilde\Theta_F$ is given by
\begin{equation}
  \label{eq:FourierRLM}
  \begin{split}
\cN_{M,L}(F,\psi) 
& = \frac{1}{L} \int_{\mathbb{R}} \hat{\psi}( L^{-1}\xi)  \bigg(\sum_{m \in \mathbb{Z}^d} F(M^{-1} m) \e\big( \tfrac{1}{2} \xi Q(m) \big) \bigg)  \dd \xi \\
&  =  \frac{M^{d-2}}{L} \int_{\mathbb{R}} \hat{\psi} ( L^{-1}\xi ) \widetilde\Theta_F\big(1, g_0 n(\xi) a(\log M) g_0^{-1}\big)  \dd \xi,
\end{split}
\end{equation}
where $\hat{\psi}$ is the Fourier transform of $\psi$,
\begin{equation}
  \label{eq:psihatdef}
  \hat{\psi}(\xi) = \int_{\mathbb{R}} \psi(x) \e(- \xi x) \dd x .
\end{equation}
The theta function is evaluated at the product of the elements
\begin{equation}
  \label{eq:g0andef}
  g_0 = \left(
  \begin{pmatrix}
    x_j^{\frac{1}{2}} & 0 \\
    0 & x_j^{-\frac{1}{2}}
  \end{pmatrix} 
    \right)_{\!\! j\leq d} \in G^d,
%    \qquad X_0 =
%  \begin{pmatrix}
%    x_1 & & \\
%        & x_2 & \\
%        & & x_3
%  \end{pmatrix}
%  ,
\end{equation}
and
\begin{equation}
  \label{eq:nadef}
  n(\xi) =
  \begin{pmatrix}
    1 & \xi \\
    0 & 1
  \end{pmatrix}
  ,
  \qquad 
   a(t) =
  \begin{pmatrix}
    \e^{-t}& 0 \\
    0 & \e^{t}
  \end{pmatrix}
  \in G,
\end{equation}
where we identify $G$ with its diagonal embedding in $G^d$.


Our analysis of the unipotent average over $\xi$ in \eqref{eq:FourierRLM} is significantly complicated by the unboundedness of the theta function.
In fact, the first term on the right of \eqref{eq:main2} comes from the contribution of $\xi$ of size at most $\tfrac{\epsilon}{M}$ for a small $\epsilon > 0$, see Section \ref{sec:smallxi}.
We view this part of the integral as analogous to a ``major arc'' in the circle method. This is analogous to the approach taken in \cite{Marklof2002,Marklof2003} and \cite{BGHM2022}.

The remaining $\xi$, the ``minor arc,'' contributes the second term in \eqref{eq:main2}, which represents the ``diagonal'' term in the two-point density.
This is obtained from the equidistribution of the translated unipotent segment.
As results on the equidistribution of unipotent flows, e.g. Ratner's theorems, typically require bounded functions, we need to control the escape of mass in \eqref{eq:FourierRLM}.
We prove the following new upper bound for the Margulis function on the space of three-dimensional lattices, which is of independent interest.
\begin{thm}
  \label{thm:escapeofmass}
  Fix $d = 3$.
  If the ratios $x_i/x_j$ are $\kappa$-Diophantine for all $i \neq j$ with $\kappa$ sufficiently close to 1,
 then for any $0 < \delta < \tfrac{1}{\kappa + 1}$, $\eta >0$ sufficiently small and $1 \leq L \leq M^{d-2} $, there is $\beta > 1$ such that
  \begin{equation}
    \label{eq:escapeofmass}
    \frac{1}{L} \int_{M^{-1-\eta} < |\xi| \leq L} \alpha_3(g_0n(\xi)a(\log M))^\beta \dd \xi =
    O_{\beta,\kappa}\bigg(1+\frac{1}{L} M^{d-2 - \delta}\bigg).
  \end{equation}
\end{thm}


Here $\alpha_d(g)$ is a Margulis function, introduced in the proof of the quantitaive Oppenheim conjecture \cite{EskinMargulisMozes1998}. It measures how high $\Gamma^d g$ is in the cusp of $\Gamma^d \backslash G^d$. As for Theorem \ref{thm:main2}, we expect the above bound to be true also for any  $d\geq 4$ under suitable Diophantine conditions.

  We note that for $L > M^{1 - \delta}$ the average (\ref{eq:escapeofmass}) is bounded.
  In comparision, \cite{BGHM2022} obtains bounds for similar (more general) averages of size roughly $M^{3\beta -2}$, which holds without Diophantine conditions.
  The stronger estimate (\ref{eq:escapeofmass}) in the special case we consider is main source for the improvements in Theorem \ref{thm:main2} over the results of \cite{BGHM2022}.
  We remark that Diophantine conditions are used in \cite{BGHM2022} to obtain pointwise bound for the Margulis functions, whereas the proof of Theorem \ref{thm:escapeofmass} uses the Diophantine condition in an essential way in controlling the average.


The proof of Theorem \ref{thm:escapeofmass} is given in Section \ref{sec:escape}.
Our strategy begins simlarly to \cite[Section~5]{EskinMargulisMozes1998}.
First a global contraction of $\alpha_3(g)^\beta$ under translates of unipotent segments (Proposition \ref{prop:globalcontraction}) is proved by combining purely local calculations with geometric properties of the cusp of $\Gamma^3 \backslash G^3$, and then \eqref{eq:escapeofmass} is proved by iterating the global contraction.

However, a straightforward adaptation of \cite{EskinMargulisMozes1998} only establishes \eqref{eq:escapeofmass} for $\beta < \tfrac{2}{3}$ while $\beta > 1$ is required for \eqref{eq:FourierRLM}.
The lack of contraction for $\beta > \tfrac{2}{3}$ is already evident in the local calculations and is related to possible repeated roots of a certain polynomial.
We regain contraction away from an exceptional set by modifying the Margulis function to account for the distance between roots of this polynomial using a strategy developed in \cite{kim2025momentsmargulisfunctionsindefinite}.
The bound \eqref{eq:escapeofmass} for $\beta > 1$ is then obtained by iterating the contraction and by controlling the contribution of the exceptional set by an avoidance estimate (Proposition \ref{prop:avoidanceMargulis}).
Our proof of this avoidance estimate uses the above-mentioned work by Lindenstrauss, Mohammadi, and Wang \cite[Theorem~1]{LMW2023a} on the polynomial effective equidistribution of one-parameter unipotent flows on $\Gamma^2 \backslash G^2$.

\section{Preliminaries}
\label{sec:prelims}

\subsection{Theta functions and Shale-Weil representation}
\label{sec:oscillatortheta}

The Heisenberg group $H_d$ is defined to be the set $\R^d \times \R^d \times \R$ with multiplication given by
\begin{equation}
  \label{eq:Hmult}
  (x_1, y_1, z_1)(x_2, y_2, z_2) 
  = (x_1 + x_2, y_1 + y_2, z_1 + z_2 - \tfrac{1}{2}(x_1\transp{y_2} - x_2\transp{y}_1)),
\end{equation}
where vectors in $\R^d$ are represented as rows.
As $\Sp(d, \R)$ preserves the symplectic form $x_1\transp y_2 - x_2\transp y_1$, $g \in G^d$ acts on $H_d$ by automorphisms
\begin{equation}
  \label{eq:gautomorphism}
  h^g = \left(
  \begin{pmatrix}
    x & y
  \end{pmatrix}
  g,
  z\right),\ \mathrm{where\ } h= (x, y, z).
\end{equation}
The group $H_d \rtimes G^d$ is then defined to have multiplication
\begin{equation}
  \label{eq:Jacobimult}
  (h_1, g_1)(h_2, g_2) = (h_1 h_2^{g_1^{-1}}, g_1g_2). 
\end{equation}

The Schr\"odinger representation $W$ of $H_d$ acts on $L^2(\R^d)$ by the unitary transformations
\begin{equation}
  \label{eq:Schrodingerdef}
  [W_d(x, y, z)f](t) = \e(-z + \tfrac{1}{2} x\transp{y}) f(t + x).
\end{equation}
For $g \in G^d$, we obtain another unitary representation $W^g$ of $H_d$ by $W^g(h) = W(h^g)$. By the Stone-von Neumann theorem on the unitary equivalence of infinite-dimensional, irreducible representations of $H_d$, there is an operator $R_d(g)$ that intertwines $W$ and $W^g$, i.e., for which
\begin{equation}
  \label{eq:Rintertwine}
  W_d^g = R_d(g)^{-1} W_d R_d(g) .
\end{equation}
These $R_d(g)$ yield the projective Shale-Weil representation of $\Sp(d,\R)$, which is central in the construction of Siegel theta functions associate with a quadratic form in $d$ variables. Since we are dealing here with diagonal forms only, it suffices to consider a subgroup of $\Sp(d, \R)$, which we identify with the $d$-power of $G=\SL(2,\R)=\Sp(1,\R)$ via the embedding
\begin{equation}
  \label{eq:gembed}
  G^d \to \Sp(d,\R),\qquad 
  \left(
  \begin{pmatrix}
    a_j & b_j \\
    c_j & d_j 
  \end{pmatrix}
  \right)_{1\leq j \leq d} \mapsto
  \begin{pmatrix}
    a_1 &  & & b_1 & &  \\
        & \ddots & & & \ddots & \\
        & & a_d & & & b_d \\
    c_1 &  & & d_1 & &  \\
        & \ddots & & & \ddots & \\
        & & c_d & & & d_d \\
  \end{pmatrix}
  .
\end{equation}
The Shale-Weil representation $R_d$ restricted to $G^d$ is then given as the $d$th tensor product of $R_1$. That is, for $g = (g_1,\ldots,g_d) \in G^d$, we have
\begin{equation}
R_d(g) = R_1(g_1) \otimes \cdots \otimes R_1(g_d) . 
\end{equation}
Proposition \ref{prop:oscillator} gives an explicit method for computing unitary operators $R_1(g)$ on a dense subset of $L^2(\R)$.

\begin{prop}
  \label{prop:oscillator}
  Let $f\in\cS(\R)$.
  For 
  \begin{equation}
    \label{eq:poscillator}
    p =
    \begin{pmatrix}
      a & b \\
      0 & a^{-1}
    \end{pmatrix}
  \end{equation}
  and any $g \in G$, we have
  \begin{equation}
    \label{eq:poscillator1}
    R_1(p)f(x) = |a|^{\frac{1}{2}} \e( \tfrac{1}{2}  ab x^2) f(ax).
  \end{equation}
  Moreover, for $ g = 
  \begin{pmatrix}
    a & b \\
    c & d
  \end{pmatrix}
  \in G$ with $c \neq 0$, we have
  \begin{equation}
    \label{eq:Fouriertransform}
    R_1(g)f(x) = |c |^{-\frac{1}{2}} \e ( \tfrac{a}{2c}x^2) \int_{\R} f(y) \e( \tfrac{d}{2c}y^2 - \tfrac{1}{c} x y) \dd y.
  \end{equation}
\end{prop}

For $f\in\cS(\R^d)$, the theta function $\Theta_f$ is defined by
\begin{equation}
  \label{eq:thetadef}
  \Theta_f(h, g) = \sum_{m \in \Z^d} [W(h)R_d(g)f](m).
\end{equation}
We observe that for $1 =(0,0,0) \in H_d$ and $g = \left(
\begin{pmatrix}
  1 & x_j \\
  0 & 1
\end{pmatrix}
\begin{pmatrix}
  y_j^{\frac{1}{2}} & 0 \\
  0 & y_j^{-\frac{1}{2}}
\end{pmatrix}
\right)_{\!\!1\leq j\leq d}$,
\begin{equation}
  \label{eq:thetaex}
  \Theta_f(1, g)
  = (y_1 \cdots y_d)^{\frac{1}{4}}\sum_{m \in \Z^d} f\left(( m_j y_j^{\frac{1}{2}})_j\right) \e\bigg( \tfrac{1}{2} \sum_{j=1}^d x_jm_j^2 \bigg).
\end{equation}
We now extend the above to a unitary representation $R_{d,d}(g)$ of $G^d$ on $L^2(\R^{2g})$
as the tensor product
\begin{equation}
R_{d,d}(g) = R_d(g) \otimes [C R_d(g) C], 
\end{equation}
where $C$ denotes complex conjugation. We highlight that $R_{d,d}$ is (unlike $R_1$, say) an actual representation for all $d$. The corresponding theta function is defined, for $F\in\cS(\R^{2d})$, by
\begin{equation}
  \label{eq:thetadefdd}
  \widetilde\Theta_F(h, g) = \sum_{m \in \Z^{2d}} [W(h)R_{d,d}(g)F](m).
\end{equation}
For example, for $F=f\otimes \overline f$ we have
$\widetilde\Theta_F(h, g) = |\Theta_f(h, g)|^2$.
Furthermore, for $g$ as in \eqref{eq:thetaex}, we have the explicit formula 
\begin{equation}
  \label{eq:thetaexdd}
  \widetilde\Theta_F(1, g)
  = (y_1 \cdots y_d)^{\frac{1}{2}}\sum_{(m,n) \in \Z^{2d}} F\left(( m_j y_j^{\frac{1}{2}})_j,( n_j y_j^{\frac{1}{2}})_j\right) \e\bigg( \tfrac{1}{2} \sum_{j=1}^d x_j(m_j^2-n_j^2) \bigg).
\end{equation}


\subsection{Automorphy of theta functions}
\label{sec:automorphy}

Define
\begin{equation}
  \label{eq:Gammabar}
  \bar{\Gamma} = \{ ((m, n, z)h_\gamma, \gamma) \subset H_d\rtimes G^d : m, n \in \Z^d,\ z \in \R, \gamma \in \Gamma^d\},
\end{equation}
where for $\gamma = (
\begin{pmatrix}
  a_j & b_j \\
  c_j & d_j
\end{pmatrix}
)_j,$ $h_\gamma = ((\tfrac{1}{2} c_jd_j)_j, (\tfrac{1}{2}a_jb_j)_j , 0)$.
We note that modulo left multiplication by $\Z^d \times \Z^d \times \R \subset H_d$,
\begin{equation}
  \label{eq:hgammamult}
  h_{\gamma_1 \gamma_2} \equiv h_{\gamma_1} (h_{\gamma_2})^{\gamma_1^{-1}},
\end{equation}
so \eqref{eq:Gammabar} does define a subgroup. The following theorem follows immediately from \cite[Theorem~4.1]{MarklofWelsh2021a}; cf.~also \cite{Marklof2003}.

\begin{thm}
  \label{thm:thetaautomorphy}
  For $F\in\cS(\R^{2d})$,  $\bar\gamma \in \bar{\Gamma}$ and $(h,g) \in H_d \rtimes G^d$,
  \begin{equation}
    \label{eq:thetaautomorphy}
     \widetilde\Theta_F\big(\bar\gamma (h, g)\big) =  \widetilde\Theta_F(h, g) .
  \end{equation}
\end{thm}

Note that in particular 
\begin{equation}
\label{thetainversion}
\widetilde\Theta_F\big((h_\gamma, \gamma)(1, g)\big) =  \widetilde\Theta_F(1, g)
\end{equation}
for all $\gamma\in\Gamma^d$.
This implies that $\widetilde\Theta_F(1, \,\cdot\,) $ defines a function on $\Gamma_{2,2}^d\backslash G^d$ with the finite index subgroup $\Gamma_{2,2}^d<\Gamma^d$ such that $c_jd_j \equiv a_jb_j \equiv 0 \bmod 2$. 

\subsection{Margulis functions}
\label{sec:Margulisfunctions}

We let $\cD$ be the standard fundamental domain for $\Gamma \backslash G$, which we take as $\cD = \cD^+ \cup \cD^-$ with
\begin{equation}
  \label{eq:standardfunddom}
  \begin{split}
    \cD^+ = & \left\{
    \begin{pmatrix}
      1 & x \\
      0 & 1
    \end{pmatrix}
    \begin{pmatrix}
      y^{\frac{1}{2}} & 0 \\
      0 & y^{-\frac{1}{2}}
    \end{pmatrix}
    k(\theta)    : 
    0 \leq x \leq \tfrac{1}{2},\ x^2 + y^2 \geq 1,\ \theta \in [0, \pi)\right\}. \\
    \cD^- = & \left\{
    \begin{pmatrix}
      1 & x \\
      0 & 1
    \end{pmatrix}
    \begin{pmatrix}
      y^{\frac{1}{2}} & 0 \\
      0 & y^{-\frac{1}{2}}
    \end{pmatrix}
    k(\theta)          : 
   -\tfrac{1}{2} < x < 0,\ x^2 + y^2 > 1,\ \theta \in [0, \pi)\right\},
  \end{split} 
\end{equation}
where
\begin{equation}
k(\theta) = \begin{pmatrix}
      \cos \theta & -\sin \theta \\
      \sin \theta & \cos \theta
    \end{pmatrix} .
\end{equation}

For $g = \left(
\begin{pmatrix}
  a_j & b_j \\
  c_j & d_j
\end{pmatrix}
\right)_{\!\!1\leq j\leq d} \in G^d$, we define the height functions
\begin{equation}
  \label{eq:rholdef}
  \rho_l(g) = \min_{\substack{J \subset \{1, \dots, d\} \\ |J| = l}} \prod_{j \in J} |(c_j, d_j)|,
\end{equation}
and
\begin{equation}
  \label{eq:alphaldef}
  \alpha_l(g) = \max_{\gamma \in \Gamma^d} \rho_l(\gamma g)^{-1}.
\end{equation}
We also use $\rho(g) := \rho_1(g)$, $\alpha(g) := \alpha_1(g)$ when $g \in G$.

For $g = (g_j)_j\in G^d$, we let $\gamma(g) = (\gamma_j(g))_j$ be so that $\gamma_j(g) g_j \in \cD$. We drop the argument and simply write $\gamma$ and $\gamma_j$ in the following.
We record the following elementary properties, which follow directly from the choice of fundemental domain $\cD$.
\begin{lem}
  \label{lem:shortestvector}
  Let $g = (g_j)_j$ and $\gamma = (\gamma_j)_j$.
  Then
  \begin{enumerate}[{\rm (i)}]
  \item $\alpha_l(g) = \min_{\substack{J \subset \{1, \dots, d\} \\ |J| = l}} \prod_{j \in J} \rho(\gamma_j g_j)$,
    
  \item $\rho(\gamma_j g_j) \ll 1$, and
    
  \item if $\rho(g_j) \leq 1$, then $\gamma_j$ has the form $
  \pm \begin{pmatrix}
    1 & * \\
    0 & 1
  \end{pmatrix}
  $.
  \end{enumerate}
\end{lem}

\begin{lem}
  \label{lem:basicfundamentaldomain}
  Set $g_0 n(\xi) a(t) = (g_j)_j$ and $\gamma( g_0 n(\xi) a(t)) = \left(
  \begin{pmatrix}
    * & * \\
    c_j & d_j
  \end{pmatrix}
 \right)_{\!\!1\leq j\leq d}$.
  \begin{enumerate}[{\rm (i)}]    
  \item
    If $\rho(g_j) \leq \delta_j$, then we have
    \begin{equation}
      \label{eq:cjdjbounds}
      |c_j| \ll \delta_j \e^t\ \mathrm{and\ } |c_j x_j \xi + d_j | \ll \delta_j \e^{-t}.
    \end{equation}
    % In particular, \eqref{eq:cjdjbounds} always holds for some $\delta_j \ll 1$.

  \item There exists $t_0 > 0$ depending only on $x_1, x_2$ so that if $t \geq t_0$, then $c_j \neq 0$.
    
  \item There exists a constant $\epsilon_0 > 0$ so that if $t \geq t_0$ and $|\xi| \geq \epsilon_0 \e^{-t}$, then also $d_1, d_2 \neq 0$. 
  \end{enumerate}
\end{lem}

The functions $\alpha_l$ play an important role in the behavior of the theta functions.
The following is a consequence of estimates for the Shale-Weil representation combined with elementary estimates for sums over lattices, see \cite[Theorem~4.4]{MarklofWelsh2021a} and \cite{Marklof2003}.

For $f\in\cS(\R^d)$, we write
  \begin{equation}
    \label{eq:ftheta}
    f_\theta = R_d\left(\left(
    \begin{pmatrix}
      \cos \theta_j & -\sin \theta_j \\
      \sin \theta_j & \cos \theta_j
    \end{pmatrix}
   \right)_{\!\!1\leq j\leq d}\right)f ,
  \end{equation}
and similarly for $F\in\cS(\R^{2d})$,
  \begin{equation}
    \label{eq:fthetadd}
    F_\theta = R_{d,d}\left(\left(
    \begin{pmatrix}
      \cos \theta_j & -\sin \theta_j \\
      \sin \theta_j & \cos \theta_j
    \end{pmatrix}
   \right)_{\!\!1\leq j\leq d}\right) F.
  \end{equation}

\begin{thm}
  \label{thm:thetaasymptotic}
  Let $F\in\cS(\R^{2d})$ and $A > 0$. Then for $h=(0,y,z)\in H_d$ and
  $$g
  = (
  \begin{pmatrix}
    1 & u_j \\
    0 & 1
  \end{pmatrix}
  \begin{pmatrix}
    v_j^{\frac{1}{2}} & 0 \\
    0 & v_j^{-\frac{1}{2}}
  \end{pmatrix}
  \begin{pmatrix}
    \cos \theta_j & -\sin \theta_j \\
    \sin \theta_j & \cos \theta_j
  \end{pmatrix}
  \in\cD,$$ we have
  \begin{equation}
    \label{eq:thetaasymp}
    \widetilde\Theta_F(h, g) = (v_1 \cdots v_d)^{\frac{1}{2}} F_\theta(0) + O_{F,A}\big((\min v_j)^{-A}\big) 
  \end{equation}
  and, for all $h\in H_d$, $g\in G^g$,
 \begin{equation}
    \label{eq:thetaasymp123}
    |\widetilde\Theta_F(h, g)|  \ll_F \alpha_d(g).
  \end{equation}
\end{thm}

In what follows we use the following estimate on the change of $\rho$ under translation by $n(\xi)$ and $a(s)$.
\begin{lem}
  \label{lem:rhotransbounds}
  Suppose that $|\xi|\leq 1$ and $s > 0$.
  Then
  \begin{equation}
    \label{eq:rhotranslationbounds}
    \tfrac{1}{2 } \e^{-s}  \rho(g) \leq \rho(g n(\xi)a(s)) \leq 2 \e^s \rho(g). 
  \end{equation}
\end{lem}

\begin{proof}
  Let $g =
  \begin{pmatrix}
    * & * \\
    c & d
  \end{pmatrix}
  ,$ so $\rho(gn(\xi)a(s)) = |
  \begin{pmatrix}
    c\e^{-s} & (c\xi + d) \e^ss
  \end{pmatrix}
  |$.
  We consider two cases.

  First, if $|c| \geq \tfrac{1}{2} |d|$, then we have
  \begin{equation}
    \label{eq:rholowerboundc}
    \rho(gn(\xi)a(s)) \geq |c| \e^{-s} \geq \tfrac{1}{2} \e^{-s} \rho(g). 
  \end{equation}
  On the other hand, if $|c| \leq \tfrac{1}{2} |d|$, then $|c\xi + d| \geq \tfrac{1}{2}|d|$ and
  \begin{equation}
    \label{eq:rholowerboundd}
    \rho(gn(\xi)a(s)) \geq \tfrac{1}{2} |d| \e^s \geq \tfrac{1}{2} \e^s\rho(g). 
  \end{equation}
  
  This establishes the lower bound in \eqref{eq:rhotranslationbounds}, and the upper bound follows by applying the lower bound with $g$ replaced with $ga(-s)n(-\xi)$. 
\end{proof}

\subsection{Diophantine property}
\label{sec:Diophantine}

\begin{lem}
  \label{lem:alpha3supbound}
  Assume that $g_0 = \left( \begin{pmatrix}   x_j^{\frac{1}{2}} & 0 \\
    0 & x_j^{-\frac{1}{2}}
  \end{pmatrix}
  \right)_{\!\! j=1,2,3}$ with $\tfrac{x_i}{x_j}$ $\kappa$-Diophantine for $i\neq j$.
  Then for any $\eta > 0$, if $\e^{-(1+\eta) t} < |\xi| \leq 10$, then $\alpha_3(g_0n(\xi) a(t)) \ll \e^{(\frac{3\kappa}{\kappa + 1} + 3\eta)t }$.
\end{lem}

\begin{proof}
  Let $g_0 n(\xi) a(t) = (g_j)_j$, $\gamma(g_0 n(\xi) a(t)) = \left(  \begin{pmatrix}  * & * \\
    c_j & d_j
  \end{pmatrix}
\right)_{\!\!j=1,2,3}$.
We may assume using Lemma \ref{lem:basicfundamentaldomain} (ii) that  $t$ is large enough so that $c_j\neq 0$.
We set $\rho(g_j) = \delta_j$, $j = 1,2,3$ and we note that $|\xi| \ll 1$ implies that $|d_j| \ll |c_j| \ll \delta_j \e^t$.

Let us suppose that $d_j \neq 0$. 
  From Lemma~\ref{lem:basicfundamentaldomain} (i) and the Diophantine condition we have
  \begin{equation}
    \label{eq:deltalowerbound1}
    (\delta_i \delta_j \e^{2t})^{-\kappa} \ll |c_i d_j \tfrac{x_i}{x_j} - c_j d_i| \ll \delta_i \delta_j,
  \end{equation}
  so $(\delta_i \delta_j)^{-1} \ll \e^{\frac{2\kappa}{\kappa + 1}t}$.
  On the other hand, if $d_j = 0$, we have
  \begin{equation}
    \label{eq:dj0lowerbound}
    \delta_j^{-1} \ll \e^{-t} + |\xi|^{-1}  \e^t \ll \e^{\eta t}.
  \end{equation}
  The claim follows.
\end{proof}

The following lemma connects the classical Diophantine notion with that used in \cite[Theorem~1]{LMW2023a}.
\begin{lem}\label{lem:periodicorbits}
    Let $L \geq 1$, $g_0=\left(\left(\begin{matrix}
        x_j^{\frac{1}{2}} & 0 \\ 0 & x_j^{-\frac{1}{2}}
    \end{matrix}\right)\right)_{\!\!j=1,2}\in G^2$ and assume that $\frac{x_1}{x_2}$ is $\kappa$-Diophantine.
Then for any $\xi\in \R$ and for any closed orbit $gH\subset  \Gamma^2\backslash G^2$, where $H=\operatorname{SL}(2, \mathbb{R})$ embedded diagonally,
    $$\operatorname{dist}_{\Gamma^2 \backslash G^2}( g_0n(\xi), gH)\gg (\operatorname{Vol}(gH))^{-D}$$
    for some constant $D = D(\kappa) >1$.
\end{lem}
\begin{proof}
  Let $$\tilde n(\xi_1, \xi_2) =
  \left(
  \begin{pmatrix}
    1 & \xi_j \\
    0 & 1
  \end{pmatrix}
\right)_{\!\! j=1,2} . $$
 We may write $\Gamma^2 g_0 n(\xi) = \Gamma^2 g_0 \tilde n(\xi_1, \xi_2)$, for suitable bounded $(\xi_1, \xi_2)\in\R^2$.
  It is therefore sufficent to show that $\operatorname{dist}_{\Gamma^2 \backslash G^2} (g_0 \tilde n(\xi_1, \xi_2), g H) \gg (\operatorname{Vol}(g H))^{-D}$
  for $|\xi_1|, |\xi_2| \ll 1$. 
  
  We may assume $g = (\mathrm{id}, h)$.
  We note that if for $A \in \R^{2 \times 2}$, $g_1 A g_2^{-1} =A$ for all $(g_1,g_2) \in g H g^{-1}$, then $A$ is a multiple of $h$.
  Since $\Gamma^2 \cap g H g^{-1}$ is Zariski dense, then the same holds if $\gamma_1 A \gamma_2^{-1} = A$ for all $(\gamma_1, \gamma_2) \in \Gamma^2 \cap g H g^{-1}$.
  Since $\gamma_1 h \gamma_2^{-1} = h$ for $(\gamma_1, \gamma_2) \in \Gamma^2 \cap g H g^{-1}$ defines a system of integer linear equations with a one-dimensional set of solutions, it follows that $h = (\det A)^{-\frac{1}{2}} A$ with $A = (a_{ij})\in \Z^{2\times 2}$.

  We may assume the entries of $A$ are coprime.  
  Since $\Gamma \cap h^{-1} \Gamma h$ is conjugate in $\Gamma$ to $\Gamma_0(\det A)$, where
  $$\Gamma_0(N) = \{
  \begin{pmatrix}
    a & b \\
    c & d
  \end{pmatrix}
  \in \Gamma : c \equiv 0 \bmod N\}
  $$ is the Hecke congruence subgroup of level $N$.
  The index of $\Gamma_0(\det A)$ in $\Gamma$ is $\geq |\det(A)|$, so $|\det A| \ll \operatorname{Vol}(gH)$.

  Suppose $\operatorname{dist}_{\Gamma^2 \backslash G^2} (g_0 \tilde n(\xi_1, \xi_2), g H)  \leq \epsilon$.
  Then there is $\gamma \in \Gamma^2$, $h_1 \in G$ with $|h_1 - \mathrm{id}| \ll \epsilon$, and $h_2 \in G$ such that $g_0 \tilde n(\xi_1, \xi_2) (\mathrm{id}, h_1)= \gamma (\mathrm{id}, h) (h_2, h_2)$.
  It follows that
  \begin{equation}
    \label{eq:x1x2diophantineapprox}
    \begin{pmatrix}
      (\tfrac{x_1}{x_2})^{-\frac{1}{2}} & (x_1 x_2)^{\frac{1}{2}} (\xi_2 - \xi_1) \\
      0 & ( \tfrac{x_1}{x_2})^{\frac{1}{2}}
    \end{pmatrix}
    = \tfrac{1}{\sqrt{\det A}} A' + O(\epsilon).
  \end{equation}
  for $A' \in \Z^{2\times 2}$.
  Inspecting the $(2,2)$-entry, the Diophantine condition on $\tfrac{x_1}{x_2}$ implies $\epsilon \gg (\det A)^{-\kappa -1}$ and the lemma follows. 
\end{proof}

Applying Lemma \ref{lem:periodicorbits} to \cite[Theorem~1]{LMW2023a}, we have the following.
\begin{prop}
  \label{prop:effectiveequidistribution}
  Let $L\geq 1$ be an integer, $g_0 = \left(  \begin{pmatrix}  x_j^{\frac{1}{2}} & 0 \\
    0 & x_j^{-\frac{1}{2}}
  \end{pmatrix}
  \right)_{\!\! j=1,2}$ with $\tfrac{x_1}{x_2}$ $\kappa$-Diophantine, and $\varphi$ be a smooth compactly supported function on $\Gamma^2 \backslash G^2$.
  Then there exists $\delta > 0$ such that
  \begin{equation}
    \label{eq:effectiveequidistribution}
    \tfrac{1}{L} \int_0^L \varphi(g_0n(\xi) a(t)) \dd \xi = \int_{\Gamma^2 \backslash G^2} \varphi(g) \dd \mu(g) + O( \cS(\varphi) \e^{-\delta t})
  \end{equation}
  where $\mu$ is the Haar probability measure on $\Gamma^2 \backslash G^2$ and $\cS(\varphi)$ is a certain Sobolev norm of $\varphi$. 
\end{prop}

\begin{proof}
  We consider the integral $\int_{0}^1 \varphi(g_0 n(l + \xi) a(t)) \dd \xi$ for $l = 0,\dots, L-1$.
  Let $R = \e^{\delta_1 t}$ for small $\delta_1 > 0$.
  By \cite[Theorem~1]{LMW2023a}, we have either
  \begin{equation}
    \label{eq:effectiveequidist1}
    \int_{0}^1 \varphi(g_0 n(l + \xi) a(t)) \dd \xi = \int_{\Gamma^2 \backslash G^2} \varphi(g) \dd \mu (g) + O(\cS(\varphi) R^{-\delta_2})
  \end{equation}
  or there is $g \in \Gamma^2 \backslash G^2$ such that $\operatorname{Vol}(g H) \leq R$ and
  \begin{equation}
    \label{eq:effectiveequidistalt}
    \operatorname{dist}_{\Gamma^2 \backslash G^2} ( g_0n(l + \xi), gH) \leq (t R)^A \e^{-t}.
  \end{equation}
  By taking $\delta_1$ to be a sufficiently small (depending only on $\kappa$), Lemma \ref{lem:periodicorbits} implies that (\ref{eq:effectiveequidistalt}) cannot hold.
  Therefore (\ref{eq:effectiveequidist1}) holds with remainder bound $\cS(\varphi) \e^{-\delta t}$ with $\delta = \delta_1 \delta_2$. 
\end{proof}


\section{Asymptotics of $\cN_{M,L}(F,\psi) $}
\label{sec:asymptotics}

In this section we determine the asymptotic behavior of $R(L,M)$ assuming Theorem~\ref{thm:escapeofmass}, which is proved in the following sections.
We recall
\begin{equation}
  \label{eq:RLMtheta}
  \cN_{M,L}(F,\psi)  = \tfrac{1}{L} \int_{\R} \hat{\psi}(\tfrac{1}{L} \xi) \widetilde\Theta_F(1, g_0 n(\xi) a(t) g_0^{-1}) \dd \xi
\end{equation}
with $t = \log M$. 


\subsection{Small $\xi$}
\label{sec:smallxi}

We apply the following with $\epsilon = \e^{-\eta t}$, $\eta > 0$ small.
\begin{prop}
Let $d\geq 1$. For $\epsilon>0$ and $A > 1$,
\begin{multline}\label{lemX1}
 \tfrac{1}{L} \int_{|\xi|\leq \epsilon \e^{-t}} \hat{\psi} ( \tfrac{1}{L}\xi ) \widetilde\Theta_F(1, g_0 n(\xi) a(t) g_0^{-1}) \dd \xi \\
= \tfrac{1}{L} \e^{(d-2)t} \hat \psi(0) \int_{\mathbb{R}^{2d}} F(v) \e( \tfrac{1}{2} \xi Q(v) ) \dd v
+ O_{F,\psi, A}(\epsilon^{A} + L^{-2} \e^{(d-4)t} + L^{-1} \epsilon^{-d + 1} \e^t) .
\end{multline}
\end{prop}

\begin{proof}
It will be convenient to write the theta function as 
\begin{equation}
\widetilde\Theta_F(1, g_0 n(\xi) a(t) g_0^{-1}) =
\widetilde\Theta_{F^0}(1, g_0 n(\xi) a(t))
\end{equation}
with $F^0 = R_{d,d}(g_0^{-1}) F$.
We set 
\begin{equation}
  \label{eq:JgXM}
  g = \left(  \begin{pmatrix}   0 & -1 \\
    1 & 0 
  \end{pmatrix}
  \right)_{\!\!1\leq j\leq d}
  g_0  n(\xi) a(t)  = \left(  \begin{pmatrix}   0 & - \e^t x_j^{-\frac{1}{2}} \\
    \e^{-t} x_j^{\frac{1}{2}} & \xi \e^t x_j^{\frac{1}{2}}
  \end{pmatrix}
  \right)_{\!\!1\leq j\leq d},
\end{equation}
and we compute that
\begin{equation}
  \label{eq:gIwasawa}
  g = \left(  \begin{pmatrix}   1 & u_j \\
    0 & 1
  \end{pmatrix}
  \begin{pmatrix}
    v_j^{\frac{1}{2}} & 0 \\
    0 & v_j^{-\frac{1}{2}}
  \end{pmatrix}
  \begin{pmatrix}
    \cos \theta_j & -\sin \theta_j \\
    \sin \theta_j  & \cos \theta_j
  \end{pmatrix}
  \right)_{\!\!1\leq j\leq d},
\end{equation}
with
\begin{equation}
  \label{eq:YRS}
  v_j =  (\e^{-2t} + \xi^2\e^{2t})^{-1}  x_j^{-1},\ \cos \theta_j = \xi \e^t x_j^{\frac{1}{2}}v_j^{\frac{1}{2}},\  \sin \theta_j = \e^{-t} x_j^{\frac{1}{2}}v_j^{\frac{1}{2}}. 
\end{equation}
We observe that if $|\xi| \leq \epsilon \e^{-t}$, then
\begin{equation}
  \label{eq:vnGammang}
  \min_{1\leq j\leq d} v_j \gg \epsilon^{-2}.
\end{equation}
By Lemma \ref{lem:shortestvector}, we may choose $\epsilon$ sufficiently small so that $\alpha_d(g) = \rho_d(g)^{-1}$ and $\alpha_{d-1}(g) \ll \epsilon \alpha_d(g)$. 

By Theorem \ref{thm:thetaasymptotic}, we have
\begin{equation}
  \label{eq:Thetaasymptotic}
  \widetilde\Theta_{F^0}( 1, g_0 n(\xi) a(t ))  = ( v_1 \cdots v_d)^{\frac{1}{2}} F^0_\theta(0) + O_{F, A}(\epsilon^A)
\end{equation}
for any $A > 0$ and from Proposition \ref{prop:oscillator} and \eqref{eq:YRS}, we have
\begin{equation}
  \label{eq:fQ0}
  F^0_\theta(0) = \e^{dt} (v_1 \cdots v_d)^{-\frac{1}{2}} \int_{\mathbb{R}^{2d}} F^0(u,u') \e( \tfrac{1}{2} \xi \e^{2t} (u\transp{u}-u'\transp{u}') ) \dd u \dd u'.
\end{equation}
It follows that
\begin{multline}
  \label{eq:smallxi}
  \tfrac{1}{L} \int_{|\xi|\leq \epsilon \e^{-t}} \hat{\psi} ( \tfrac{1}{L}\xi ) \widetilde\Theta_{F^0}(1, g_0 n(\xi) a(t) ) \dd \xi \\
  = \tfrac{1}{L}\e^{dt} \int_{|\xi| \leq \epsilon \e^{-t}} \hat{\psi}( \tfrac{1}{L}\xi)  \bigg( \int_{\mathbb{R}^{2d}} F^0(u,u') \e( \tfrac{1}{2} \xi \e^{2t} (u\transp{u}-u'\transp{u}') ) \dd u \dd u'\bigg)  \dd \xi + O_{f, \psi, A}(\epsilon^A).
\end{multline}
The inner integral over $\dd u \dd u'$ is $O_{F}( (1 + |\xi| \e^{2t})^{-d})$, so replacing $\hat \psi(\tfrac{1}{L} \xi)$ with $\hat \psi(0)$ creates an error bounded by
\begin{equation}
  \label{eq:psi0replace}
  \ll_{\psi, F} L^{-2} \e^{dt} \int_{|\xi| \leq \epsilon \e^{-t}} \frac{ |\xi|}{(1 + |\xi| \e^{2t})^d} \dd \xi \ll L^{-2} \e^{(d-4)t}. 
\end{equation}
After changing variables $\xi \gets \e^{-2t} \xi$, we have that (\ref{eq:smallxi}) is 
\begin{equation}
\begin{split}
  \label{eq:smallxi1}
&\tfrac{1}{L} \e^{(d-2)t} \hat\psi(0)  \int_{|\xi| \leq \epsilon \e^{t}} \bigg( \int_{\mathbb{R}^{2d}} F^0(u,u') \e( \tfrac{1}{2} \xi (u\transp{u}-u'\transp{u}') ) \dd u \dd u'\bigg)  \dd \xi \\
&  \qquad + O_{F,\psi,A}(\epsilon^{A} + L^{-2} \e^{(d-4)t})\\
& = \tfrac{1}{L} \e^{(d-2)t} \hat \psi(0)  \int_\R \bigg( \int_{\mathbb{R}^{2d}} F^0(u,u') \e( \tfrac{1}{2} \xi (u\transp{u}-u'\transp{u}') ) \dd u \dd u'\bigg)  \dd \xi \\
 & \qquad + O_{F,\psi, A}(\epsilon^{A} + L^{-2} \e^{(d-4)t} + L^{-1} \epsilon^{-d + 1} \e^{-t}).
\end{split}
\end{equation}
Finally, we recall that $F_0=R_{d,d}(g_0^{-1}) F$, which yields after the appropriate variable substitution
\begin{equation}
\begin{split}
 & \int_{\mathbb{R}^{2d}} F^0(u,u') \e( \tfrac{1}{2} \xi (u\transp{u}-u'\transp{u}') ) \dd u \dd u' \\
& = (\det X_0)^{-1} \int_{\mathbb{R}^{2d}} F(u X_0^{-1/2},u' X_0^{-1/2}) \e( \tfrac{1}{2} \xi (u\transp{u}-u'\transp{u}') ) \dd u \dd u' \\
& = \int_{\mathbb{R}^{2d}} F(v) \e( \tfrac{1}{2} \xi Q(v) ) \dd v.
\end{split}
\end{equation}
 
\end{proof}


\subsection{Remaining $\xi$}
\label{sec:remaining}

To analyze the integral \eqref{eq:RLMtheta} over $|\xi| \geq \e^{-(1 + \eta)t}$, we appeal to results on the equidistribution of unipotent flows, ie Ratner's theorems and their relatives.

We let $\Gamma' \subset \Gamma^d $ be the finite index subroup of $\gamma = \left(
\begin{pmatrix}
  a_j & b_j \\
  c_j & d_j
\end{pmatrix}
\right)_{\!\!1\leq j\leq d}$ such that $c_jd_j \equiv a_jb_j \equiv 0 \bmod 2$.
We note that from Theorem~\ref{thm:thetaautomorphy}, $|\Theta_f(1, g)|^2$ is a well-defined function $\Gamma' \backslash G^d$.

\begin{lem}
  \label{lem:density}
  Let $x = (x_j)_j \in \R^d$ be so that $q\transp x \neq 0$ for any nonzero $q \in \Z^d$ and let $g_0 = \left(  \begin{pmatrix}   x_j^{\frac{1}{2}} & 0 \\
    0 & x_j^{-\frac{1}{2}}
  \end{pmatrix}
  \right)_{\!\!1\leq j\leq d} \in G^d$.
  Then $\Gamma' g_0 H g_0^{-1}$ is dense in $\Gamma' \backslash G^d$.
\end{lem}

\begin{proof}
  For $\xi = (\xi_j)_j \in \R^d$, we let $\tilde n(\xi) = \left(  \begin{pmatrix}   1 & \xi_j \\
    0 & 1
  \end{pmatrix}
  \right)_{\!\!1\leq j\leq d} \in G^d$.
  We have $g_0 n(\xi)a(t) g_0^{-1} = \tilde n(\xi x) a(t)$.
  Moreover, by the condition on $x$, we have that $\xi x$ is dense in $\R^d / 2\Z^d$.

  Since $\tilde n(\xi)$ for $\xi \in \R^d / 2\Z^d$ parametrizes the full horospherical subgroup for $a(t)$, it follows that the set $\{\Gamma' \tilde n(\xi)a(t) : \xi \in \R^d / 2Z^d\}$ is dense in in $\Gamma' \backslash G^d$.
  It follows that $\{\Gamma' g_0 n(\xi) a(t) g_0^{-1} : \xi \in \R\}$ is dense in $\Gamma' \backslash G^d$.
\end{proof}

By a theorem of Shah \cite{MR1403756}, Lemma \ref{lem:density} implies that for any bounded continuous function $\varphi$ on $\Gamma' \backslash G^d$,
\begin{equation}
  \label{eq:equidistribution}
  \lim_{t\to \infty} \tfrac{1}{L} \int \hat\psi( \tfrac{\xi}{L}) \varphi(g_0n(\xi)a(t)) \dd \xi = \psi(0)\int_{\Gamma' \backslash G^d} \varphi(g) \dd \mu(g),
\end{equation}
where $\mu$ is the $G^d$-invariant probability measure on $\Gamma' \backslash G^d$.
We now use Theorem~\ref{thm:escapeofmass} (proved below) to extend \eqref{eq:equidistribution} to the unbounded theta function.
We specialize to $d = 3$.

We let $\cC_y = \{g \in \Gamma^3 \backslash G^3 : \alpha_3(g) \geq y\}$, and we let $\chi_y$ be a continuous function on $\Gamma' \backslash G^3$ satisfying $\mathbbm{1}_{\cC_{2y}} \leq \chi_y \leq \mathbbm{1}_{\cC_y}$.
By theorems \ref{thm:thetaasymptotic} and \ref{thm:escapeofmass}, we have
\begin{multline}
  \label{eq:tailestimate}
  \tfrac{1}{L} \int_{|\xi| \geq \e^{-(1 + \eta)t}} |\widetilde\Theta_F(1, g_0n(\xi) a(t))| \, \chi_y(g_0n(\xi)g(t)) \dd \xi \\
  \leq \tfrac{1}{L} y^{1 - \beta} \int_{|\xi| \geq \epsilon \e^{-t}} \alpha_d(g_0n(\xi) a(t))^\beta \dd \xi \ll y^{1-\beta}( 1 + \tfrac{1}{L} M^{(3\beta -2) (\frac{\kappa}{\kappa + 1} + \eta)})
\end{multline}
for some $\beta > 1$.

On the other hand, $|\Theta(h, g)|^2(1 - \chi_y(g)) \ll y$, so \eqref{eq:equidistribution} implies that
\begin{equation}
\begin{split}
  \label{eq:equidistribution1}
&  \tfrac{1}{L} \int_{\e^{-(1 + \eta) t} \leq |\xi| \leq L} \psi(\tfrac{1}{L} \xi) \widetilde\Theta_F(1, g_0n(\xi)a(t)) (1 - \chi_y(g_0n(\xi)a(t)))\dd \xi \\
&  = \tfrac{1}{L} \int_{|\xi| \leq L} \psi(\tfrac{1}{L} \xi) \widetilde\Theta_F(1, g_0n(\xi)a(t)) (1 - \chi_y(g_0n(\xi)a(t)))\dd \xi + O_\psi(\tfrac{1}{L}y\e^{-t})\\
&  \underset{t\to\infty}{\to} \psi(0) \int_{\Gamma' \backslash G^d} \widetilde\Theta_F(1, g) ( 1 - \chi_y(g)) \dd \mu(g) \\
&  = \psi(0) \int_{\Gamma_d \backslash G^d} \widetilde\Theta_F(1, g) \dd \mu(g) + O(y^{-1}),
\end{split}
\end{equation}
where the last line follows from $\mu(\cC(y) \setminus \cC(2y)) \ll y^{-2}$. 

Combining this with \eqref{eq:tailestimate}, we obtain
\begin{equation}
  \label{eq:ratner1}
  \lim_{t \to \infty} \tfrac{1}{L} \int_{\epsilon \e^{-t} \leq |\xi| \leq L} \psi(\tfrac{1}{L} \xi) \widetilde\Theta_F(1, g_0n(\xi)a(t)) \dd \xi \\
  = \psi(0) \int_{\Gamma' \backslash G^d} \widetilde\Theta_F(1, g) \dd\mu(g). 
\end{equation}
Together with Lemma \ref{lem:thetaint} below, this gives the second term in \eqref{eq:main2}.

\begin{lem}
  \label{lem:thetaint}
  For $F\in \cS(\R^{2d})$, $d\geq 1$,
  \begin{equation}
    \label{eq:thetaint}
    \int_{\Gamma' \backslash G^d} \widetilde\Theta_F(1, g) \dd \mu(g) = \sum_{\sigma\in D_{2h}} \int_{\R^d} F(x,\sigma(x)) \dd x.
  \end{equation}
\end{lem}

\begin{proof}
  As $\tilde n$ parametrizes the full horospherical subgroup associated to $a(t)$, it follows from the mixing property of $a(t)$ via standard arguments (e.g. Margulis' thickening trick) that
  \begin{equation}
    \label{eq:fullhorosphere}
    \int_{\Gamma' \backslash G^d} \widetilde\Theta_F(1, g) \dd \mu(g) = \lim_{t\to \infty} 2^{-d} \int_{[0,2]^d} \Theta_F(1, \tilde n(\xi)a(t)) \dd \xi.
  \end{equation}
  From \eqref{eq:thetaex} we have
  \begin{equation}
  \begin{split}
    \label{eq:horosphereaverage}
    & 2^{-d} \int_{[0,2]^d} \widetilde\Theta_F(1, \tilde n(\xi)a(t)) \dd \xi \\
    & = 2^{d}\e^{-dt}\sum_{m_1, m_2 \in \Z^d} F(m_1\e^{-t}, m_2e^{-t}) \int_{[0,2]^d} \e(\tfrac{1}{2} \sum_{1\leq j\leq d} \xi_j (m_{1,j}^2 - m_{2,j}^2)) \dd \xi \\
    & = \e^{-dt}\sum_{m \in \Z^d}\sum_{\sigma\in D_{2h}} F(m\e^{-t},\sigma(m) e^{-t}).
  \end{split}
  \end{equation}
  As $t\to \infty$ this Riemann sum converges to the right side of \eqref{eq:thetaint}.
\end{proof}

\section{Escape of mass}
\label{sec:escape}

In this section we prove Theorem \ref{thm:escapeofmass} assuming the avoindance estimate (Proposition~\ref{prop:avoidanceMargulis}) proved in Section \ref{sec:avoidance}.


\subsection{Local contraction}
\label{sec:loccontraction}

The main result of this section is Lemma \ref{lem:phicontraction}, which establishes a local contraction property of the height functions $\rho_k$, compare to \cite[Lemma~5.1]{KleinbockMargulis1999}.
In our setting, we modify these height functions to also measure the distance to an exceptional set on which $\rho_k^{-\beta}$ does not contract for $\beta > \tfrac{2}{k}$.
In this way, we obtain contraction for the modified functions for $1 < \beta < 2$.

The distance to the exceptional set is controlled by the following functions.
For $g =
\left(
\begin{pmatrix}
  a_j & b_j \\
  c_j & d_j
\end{pmatrix}
\right)_{\!\! j=1,2,3} \in G^3$ we define
\begin{equation}
  \label{eq:Delta2def}
  \Delta_2(g) = \min_{j = 1,2,3} \max_{i \neq j} \min ( 1, | \tfrac{d_i}{c_i} - \tfrac{d_j}{c_j}|)
\end{equation}
and
\begin{equation}
  \label{eq:Deltadef}
  \Delta_3(g) = \min_{\!\! j=1,2,3} \prod_{i \neq j} \min(1, |\tfrac{d_i}{c_i} - \tfrac{d_j}{c_j}|),
\end{equation}
where here and in what follows, if $c_ic_j = 0$, we take $1$ in corresponding minimum.


\begin{lem}
  \label{lem:measbound}
  Let $g = \left( \begin{pmatrix}    r_j & -s_j \\
    s_j & r_j
  \end{pmatrix}
  \right)_{\!\! j=1,2,3} \in \SO(2)^3$.
  For any $s > 0$, the measure of the set of $\xi \in [-1,1]$ such that $\rho_3(gn(\xi)a(s)) \leq \delta$ is $O(\min( \delta^{\frac{1}{3}}\e^{-s} , \delta \Delta_2 (g)^{-1} \e^{-s}, \delta \Delta_3(g)^{-1} \e^{-3s}))$. 
\end{lem}

\begin{proof}
  We note that since $r_j^2 + s_j^2 = 1$, if $s_j^2 \leq \tfrac{1}{10}$, we have $|r_j + s_j \xi| \asymp 1$ for $\xi \in [-1,1]$.
  Therefore
  \begin{equation}
  \begin{split}
    \label{eq:rhokgna}
    \rho_3 (gn(\xi) a(s)) & = \prod_{j = 1}^3 |(\e^{-s} s_j, \e^s(r_j + s_j \xi))| \\
    &\gg \e^{3s} \prod_{s_j^2 > \tfrac{1}{10}} |r_j + s_j \xi| \gg \e^{3s} \prod_{s_j^2 > \tfrac{1}{10}} |\xi + \tfrac{r_j}{s_j}|,
  \end{split}
  \end{equation}
  and we may assume that product is nonempty since otherwise the lemma is trivial.

  Let $j$ be so that $|\xi + \tfrac{r_j}{s_j}|$ is minimal among $j$ so that $s_j^2 > \tfrac{1}{10}$.
  If $|\tfrac{r_j}{s_j}| > 2$, then the lemma is again trivial.
  It follows that (\ref{eq:rhokgna}) is
  \begin{equation}
    \label{eq:rho3cubelowerbound}
    \gg \e^{3s} |\xi + \tfrac{r_j}{s_j}|^3. 
  \end{equation}
  We also have that (\ref{eq:rhokgna}) is
  \begin{equation}
    \label{eq:rho3linearlowerbound}
    \gg \e^{3s} |\xi + \tfrac{r_j}{s_j}| \prod_{\substack{s_i^2 > \tfrac{1}{10} \\ i \neq j}} | \tfrac{r_i}{s_i} - \tfrac{r_j}{s_j} | \gg \e^{3s} \Delta_3(g) |\xi + \tfrac{r_j}{s_j}|, 
  \end{equation}
  where the last bound follows from $| \tfrac{r_i}{s_i} - \tfrac{r_j}{s_j} |\gg 1$ if $s_i^2 \leq \tfrac{1}{10}$.
  
  The first and third bounds claimed in the lemma follow from (\ref{eq:rho3cubelowerbound}) and (\ref{eq:rho3linearlowerbound}).
  For the second claimed bound, we aim to show
  \begin{equation}
    \label{eq:rho3quadlowerbound1}
    \rho_3(gn(\xi)a(s)) \gg \e^s \Delta_2(g) |\xi + \tfrac{r_j}{s_j}|
  \end{equation}
  with $j$ be as before.  

  Without loss of generality we assume $j = 3$.
  If both $s_1^2 \leq \tfrac{1}{10}$ and $s_2^2 \leq \tfrac{1}{10}$, then $\Delta_2(g) = \Delta_3(g) = 1$, so (\ref{eq:rho3quadlowerbound1}) follows from (\ref{eq:rho3linearlowerbound}).
  If $s_1^2 \leq \tfrac{1}{10} < s_2^2$, then from the left of (\ref{eq:rhokgna}) we have
  \begin{equation}
    \label{eq:rho3quadlowerbound4}
    \rho_3(gn(\xi)a(s)) \geq \e^s |s_2| |r_1 + s_1 \xi| |r_3 + s_3 \xi| \gg \e^s |\xi + \tfrac{r_3}{s_3}|,
  \end{equation}
  and (\ref{eq:rho3quadlowerbound1}) follows similarly if $s_2^2 \leq \tfrac{1}{10} < s_1^2$.

  Finally, if both $s_1^2, s_2^2 \leq \tfrac{1}{10}$, we assume without loss of generality that $|\tfrac{r_1}{s_1} - \tfrac{r_3}{s_3}| \geq |\tfrac{r_2}{s_2} - \tfrac{r_3}{s_3}|$.
  We have
  \begin{equation}
    \label{eq:rho3quadlowerbound5}
    \rho_3(gn(\xi)a(s)) \geq \e^s |s_2| |r_1 + s_1 \xi| |r_3 + s_3 \xi| \gg \e^s |\tfrac{r_1}{s_1} - \tfrac{r_3}{s_3}||\xi + \tfrac{r_3}{s_3}|,
  \end{equation}
  so (\ref{eq:rho3quadlowerbound1}) follows. 
\end{proof}

\begin{lem}
  \label{lem:rhocontraction}
  For any $s > 0$, $\beta > 1$, and $g \in G$, we have
  \begin{equation}
    \label{eq:loccontraction}
    \int_{-1}^1 \rho_3 (g n(\xi) a(s))^{-\beta} \dd \xi 
    \ll \min (\e^{(3\beta -2)s}, \Delta_2(g)^{-1} \e^{(3\beta - 4)s}, \Delta_3(g)^{-1} \e^{(3\beta -6)s}) \rho_k(g)^{-\beta}.
  \end{equation}
\end{lem}

\begin{proof}
  We note if $g' = \left(
  \begin{pmatrix}
    * & * \\
    0 & * 
  \end{pmatrix}\right)_{\!\! j=1,2,3}$, then $\rho_3(g'g) = \rho_3(g')\rho_3(g)$, so we may assume $g =
  \left(
  \begin{pmatrix}
    r_j & -s_j \\
    s_j & r_j
  \end{pmatrix}\right)_{\!\! j=1,2,3}
  \in \SO(2)^3$, where $\rho_3(g) = 1$.
  We note that since $r_j^2 + s_j^2 = 1$, it follows that if $s_j^2 \leq \tfrac{1}{10}$, then $|r_j + s_j \xi| \gg 1$ for all $\xi \in [-1,1]$.
  It follows that $\rho_3(gn(\xi)a(s)) \gg \e^{-(3 - 2k)s} \geq \e^{-3s}$, where $k$ is the number of $j$ such that $s_j^2 \leq \tfrac{1}{10}$.
  We also note that $\rho_3(gn(\xi)a(s)) \ll \e^{3s}$. 

  We decompose the region of integration according to $\xi$ such that $\delta < \rho_3(gn(\xi) a(s)) \leq 2 \delta$ ranges dyadically from $\e^{-3s}$ to $\e^{3s}$ ($\delta = c 2^j \e^{-3s}$, $0 \leq j \ll s$, and $c > 0$ a small constant).
  We obtain
  \begin{equation}
  \begin{split}
    \label{eq:dyadicdecomp}
    \int_{-1}^1 \rho_3 (g n(\xi) a(s))^{-\beta} \dd \xi 
&    \ll \sum_{\substack{ \delta \gg \e^{-3s} \\ \mathrm{dyadic}}} \delta^{-\beta} \mathrm{meas}\{ \xi \in [-1,1]: \rho_3(gn(\xi)a(s)) \leq \delta\} \\
 &   \ll \sum_{\substack{ \delta \gg \e^{-3s} \\ \mathrm{dyadic}}} \delta^{-\beta} \min( \delta^{\frac{1}{3}} \e^{-s}, \delta \Delta_2(g)^{-1} \e^{-s}, \delta \Delta_3^{-1} \e^{-3s}) \\
 &   \ll \min (\e^{(3\beta -2)s}, \Delta_2(g)^{-1} \e^{(3\beta - 4)s}, \Delta_3(g)^{-1} \e^{(3\beta -6)s}),
  \end{split}
  \end{equation}
  the second bound following from Lemma \ref{lem:measbound}.
\end{proof}

We set
\begin{equation}
  \label{eq:phidef}
  \phi(g) = \Delta_3(g)^{-1} \rho_3(g)^{-\beta}.
\end{equation}
\begin{lem}
  \label{lem:phicontraction}
  For any $g \in G$, $s > 0$, and $\beta > 1$, we have
  \begin{equation}
    \label{eq:phicontraction}
    \int_{-1}^1 \phi(gn(\xi)a(s)) \dd \xi \ll \e^{-3 (2-\beta) s} \phi(g).
  \end{equation}
\end{lem}

\begin{proof}
  We observe that 
  \begin{equation}
    \label{eq:kappaas}
    \Delta_3(gn(\xi)a(s))^{-1} \Delta_3(g) =
    \begin{cases}
      \e^{-4s}  & \mathrm{if\ } \Delta_2(g) \leq \e^{-2s} \\
      \e^{-2s} \Delta_2(g)  & \mathrm{if\ }  \Delta_3(g) \leq \e^{-2s} \leq \Delta_2(g) \\
      \Delta_3(g) & \mathrm{if\ } \e^{-2s} \leq \Delta_3(g).
    \end{cases}
  \end{equation}
  Applying the bounds in Lemma \ref{lem:rhocontraction} to each of these cases (in the corresponding order) implies (\ref{eq:phicontraction}). 
\end{proof}

We end this section by recording the following lemma establishing a weaker bound than Lemma~\ref{lem:phicontraction}, but for a larger range of $\beta$ that we omitted from the proof of Lemma~\ref{lem:phicontraction} for the sake of exposition.
\begin{lem}
  \label{lem:cruderhocontraction}
  For any $d\geq 1$, $1\leq k\leq d$, $s > 0$, $\beta > \tfrac{1}{k}$, and $g \in G^d$, we have
  \begin{equation}
    \label{eq:crudeloccontraction}
    \int_{-1}^1 \rho_k (g n(\xi) a(s))^{-\beta} \dd \xi \ll \e^{-(2-\beta k )s} \rho_k(g)^{-\beta}.
  \end{equation}
\end{lem}

\begin{proof}
  As in the proof of Lemma \ref{lem:rhocontraction}, for $g' = \left(\begin{pmatrix}    * & * \\
    0 & * 
  \end{pmatrix}
  \right)_{\!\!1\leq j\leq d}$, we have $\rho_k(g'g) = \rho_k(g')\rho_k(g)$, so we may assume that $g = \left(\begin{pmatrix}    r_j & -s_j \\
    s_j & r_j
  \end{pmatrix}
  \right)_{\!\!1\leq j\leq d} \in \SO(2)^d$, where we have $\rho_k(g) = 1$.

  We have
  \begin{equation}
    \label{eq:deltalowerbound}
    \rho_k(g) \asymp \min_{\substack{J \subset\{1,\dots, d\} \\ |J|=k}} \max_{J_1 \subset J} \e^{-(k-2|J_1|)s}\prod_{j \in J \setminus J_1} |s_j| \prod_{j \in J_1} |r_j + s_j \xi|.
  \end{equation}
  Let $J_0$ be the set of $j$ such that $s_j^2 \leq \tfrac{1}{10}$, so that if $j \in J_0$, $|r_j + s_j \xi| \asymp 1$ for $\xi \in [-1,1]$.
  Taking $J_1 = J_0 \cap J$ in (\ref{eq:deltalowerbound}) shows that $\rho_k(gn(\xi)a(s)) \gg \e^{-(k-|J_0 \cap J|)s} \geq \e^{-ks}$.

  On the other hand, we claim that the measure of $\xi \in [-1,1]$ such that $\rho_k(gn(\xi)a(s)) \leq \delta $ is $\ll \delta^{\frac{1}{k}} \e^{-s}$.
  Taking $J_1 = J$, we have
  \begin{equation}
    \label{eq:rhoklowerbound}
    \rho_k(gn(\xi)a(s)) \gg \e^{ks} \prod_{j \in J \setminus J_0} |r_j + s_j \xi| \gg \e^{ks} \prod_{j \in J \setminus J_0} |\xi + \tfrac{r_j}{s_j}|. 
  \end{equation}
  Let $j \in J\setminus J_0$ be so that $|\xi + \tfrac{r_j}{s_j}|$. If $|\tfrac{r_j}{s_j}| \geq 2$ then the claim is trivial.
  Otherwise we have
  $$\rho_k(gn(\xi)a(s)) \gg \e^{ks} |\xi + \tfrac{r_j}{s_j}|^{|J \setminus J_0|} \gg \e^{ks} |\xi + \tfrac{r_j}{s_j}|^{k},$$
  and the claim follows in this case as well.

  Decomposing $\xi \in [-1,1]$ according to $\tfrac{1}{2} \delta < \rho_k(gn(\xi)a(s)) \leq \delta $, it follows that
 \begin{equation}
  \begin{split}
    \label{eq:dyadicdecomp2}
    \int_{-1}^1 \rho_k (g n(\xi) a(s))^{-\beta} \dd \xi 
& \ll \sum_{\substack{ \delta \gg \e^{-ks} \\ \mathrm{dyadic}}} \delta^{-\beta} \mathrm{meas}\{ \xi \in [-1,1]: \rho_k(gn(\xi)a(s)) \leq \delta\} \\
 &   \ll \sum_{\substack{ \delta \gg \e^{-ks} \\ \mathrm{dyadic}}} \delta^{\frac{1}{k}-\beta}\e^{-s}  \ll \e^{-(2 - \beta k)s},
  \end{split}
  \end{equation}
  as required. 
\end{proof}

\subsection{Global contraction}
\label{sec:globalcontraction}

We first prove a crude bound for averages of $\alpha_d(gn(\xi)a(s))^{\beta}$.
We note that this bound does not exhibit contraction for $\beta \geq \tfrac{2}{d}$.
\begin{lem}
  \label{lem:crudecontraction}
  Let $d\geq 1$. For any $g \in G^d$, $s > 0$, and $\beta > \tfrac{1}{d}$, we have
  \begin{equation}
    \label{eq:crudecontraction}
    \int_{-1}^1 \alpha_d(gn(\xi)a(s))^\beta \dd \xi 
    \ll \e^{(\beta d- 2)s} \alpha_d(g)^\beta + \e^{\beta s}\int_{-1}^1 \alpha_{d-1}(gn(\xi)a(s))^\beta d \xi .
  \end{equation}
  In particular, we have
  \begin{equation}
    \label{eq:crudecontraction1}
    \int_{-1}^1 \alpha_d(gn(\xi)a(s))^\beta \dd \xi \ll \e^{(\beta d- 2)s} \alpha_d(g)^\beta + \e^{\beta d s}.
  \end{equation}
\end{lem}

\begin{proof}
  We let $\gamma(\xi) = \gamma(gn(\xi) a(s))$ and we note that \eqref{eq:crudecontraction} follows from the following
  \begin{equation}
    \label{eq:pointwisecrudecontraction}
    \rho_d(gn(\xi)a(s))^\beta \dd \xi \ll \rho_d(\gamma(g) g)^\beta + \e^{2s}\alpha_{d-1}(gn(\xi)a(s))^\beta
  \end{equation}
  together with Lemma \ref{lem:cruderhocontraction}. 

  If $\gamma(\xi) = \left(\begin{pmatrix}    \pm 1 & * \\
    0 & \pm 1
  \end{pmatrix}
  \right)_{\!\!1\leq j\leq d} \gamma(g)$, then \eqref{eq:pointwisecrudecontraction} is immediate.
  If for some $j$, $\gamma_j(\xi) \neq \left(\begin{pmatrix}    \pm 1 & * \\
    0 & \pm 1
  \end{pmatrix}
  \right)_{\!\!1\leq j\leq d} \gamma_j(g)$, then it follows from Lemma \ref{lem:shortestvector} that $\rho(\gamma_j(\xi)g_j) > 1$, so by Lemma \ref{lem:rhotransbounds}
  \begin{equation}
    \label{eq:rhogammajxichange}
    \rho(\gamma_j(\xi)g_j n(\xi) a(s)) \geq \tfrac{1}{2} \e^{-s} \rho(\gamma_j(\xi)g_j) \geq \tfrac{1}{2} \e^{-s}.
  \end{equation}
  We now have
  \begin{equation}
    \label{eq:gammachangebound}
    \rho_d(\gamma(\xi)gn(\xi)a(s)) \leq 2 \e^s \alpha_{d-1}(gn(\xi)a(s))
  \end{equation}
  and \eqref{eq:pointwisecrudecontraction} follows in this case as well.

  Finally, we observe that \eqref{eq:crudecontraction1} follows from \eqref{eq:crudecontraction} inductively. 
\end{proof}

In order to obtain a contraction, we define a global modified height function $\widetilde{\alpha}: \Gamma^3 \backslash G^3 \to (0,\infty)$ by
\begin{equation}
  \label{eq:tildealphadef1}
  \widetilde{\alpha}(g) = \phi(\gamma(g)g).
\end{equation}
We now consider exceptional sets to the global contraction hypothesis.
For $s>0$ and $K>1$ we define
\begin{multline}
  \label{eq:cEdef}
  \mathcal{E}_{s,K}:=\bigg\{\Gamma^3 g\in \Gamma^3 \backslash G^3: \Delta_3(\gamma(gn(\xi)a(s)) g)<\e^{-Ks} \\
  \textrm{ for some } 
  \xi \in [-1,1] \textrm{ such that }  \gamma(gn(\xi)a(s)) \neq \left(\begin{pmatrix}    \pm 1 & * \\
    0 & \pm 1
  \end{pmatrix}
  \right)_{\!\! j=1,2,3} \gamma(g)
  \bigg\}.
\end{multline}
The following global contraction hypothesis holds outside the set $\mathcal{E}_{s,K}$.
\begin{prop}
  \label{prop:globalcontraction}
  Let $1<\beta<2$.
  For any $s>0$, $K>1$, and $\Gamma^3 g \notin \mathcal{E}_{s,K}$ we have
  \begin{equation}
    \label{eq:globalcontraction}
    \int_{-1}^1\widetilde{\alpha}\big(gn(\xi)a(s)\big)\dd\xi \\
    \ll  e^{-3(2-\beta) t}\widetilde{\alpha}(g)+e^{(K+2)s}\int_{-1}^{1} \alpha_{2}(gn(\xi)a(s))^{\beta}\dd \xi.
  \end{equation}
\end{prop}

\begin{proof}
  We argue as in the proof of Lemma \ref{lem:crudecontraction}.
  It suffices to show that
  \begin{equation}
    \label{eq:pointwisecontractionbound}
    \phi(\gamma(\xi)gn(\xi)a(s)) \ll \phi(\gamma(g)gn(\xi)a(s)) +e^{(K+2)s} \alpha_{2}(gn(\xi)a(s))^{\beta},
  \end{equation}
  since then \eqref{eq:globalcontraction} follows from Lemma \ref{lem:phicontraction}.

  If $\gamma(\xi) = \left(\begin{pmatrix}    \pm 1 & * \\
    0 & \pm 1
  \end{pmatrix}
  \right)_{\!\! j=1,2,3} \gamma(g)$, then \eqref{eq:pointwisecontractionbound} is immediate, so let us suppose that there is $j$ such that $\gamma_j(\xi) \neq \left(\begin{pmatrix}    \pm 1 & * \\
    0 & \pm 1
  \end{pmatrix}
  \right)_{\!\! j=1,2,3} \gamma_j(g)$.
  Then we have $\rho(\gamma_j(\xi)g_jn(\xi)a(s)) \geq \tfrac{1}{2} \e^{-s}$ and therefore
  \begin{equation}
    \label{eq:gammachangebound1}
    \alpha_3(gn(\xi) a(s)) \leq 2\e^s \alpha_{2}(gn(\xi)a(s))
  \end{equation}
  as in the proof of Lemma \ref{lem:crudecontraction}.
  Since $g \not\in \cE_{s,K}$, \eqref{eq:pointwisecontractionbound} follows.
\end{proof}

In what follows it is at times easier to work with the following set.
  \begin{lem}
    \label{lem:altexceptionalset}
    We have
    \begin{multline}
      \label{eq:tildecEdef}
      \cE_{s,K} \subset \tilde{\cE}_{s,K} := \{\Gamma^3 g\in \Gamma^3 \backslash G^3: \Delta_3(\gamma g)<\e^{-Ks} \\
  \textrm{ for some } 
  \gamma \in \Gamma^3
  \textrm{ such that } \tfrac{1}{10}\e^{-s} < \max_{k=1,2,3}\rho(\gamma_kg_k) \leq 10 \e^{s} \}.
    \end{multline}
  \end{lem}
  
\begin{proof}
  As above,
  $\gamma_j(\xi) \neq \left(\begin{pmatrix}    \pm 1 & * \\
    0 & \pm 1
  \end{pmatrix}
  \right)_{\!\! j=1,2,3} \gamma_j(g)$ implies $\rho(\gamma_j(\xi)g_jn(\xi)a(s)) \geq \tfrac{1}{2} \e^{-s}$.
  We also have from Lemma \ref{lem:rhotransbounds} that
  \begin{equation}
    \label{eq:rhoupperbound}
    \rho(\gamma_j(\xi)g_j) \leq 2\e^s \rho(\gamma_jn(\xi)a(s)) \leq 10 \e^s,
  \end{equation}
  so the containment \eqref{eq:tildecEdef} follows.
\end{proof}

\subsection{Smaller moments}
\label{sec:smallmoments}

In this section we prove Proposition \ref{prop:smallmoment}, which controls the escape of mass for $\alpha_3^\beta$ with $\beta < \tfrac{2}{3}$.
Combined with Lemma \ref{lem:alpha3supbound} this always for some useful preliminary control over the escape of mass.
\begin{prop}
  \label{prop:smallmoment}
  Let $d\geq 1$. For any $\Gamma^dg \in \Gamma^d \backslash G^d$ and $\beta < \tfrac{2}{d}$, we have
  \begin{equation}
    \label{eq:smallmoment}
    \sup_{t > 0} \int_{-10}^{10} \alpha_d(gn(\xi) a(t))^\beta \dd \xi < \infty.
  \end{equation}
\end{prop}

\begin{proof}
  We prove this by induction on $d$.
  The case $d = 1$ being verified by the following straightforward calculation for $\beta < 2$:
  \begin{multline}
    \label{eq:d1smallmoment}
    \int_{-10}^{10} \alpha_1(gn(\xi) a(t))^\beta \dd \xi \ll \int_{|\xi| \leq \epsilon \e^{-t}} ( \e^{-2t} + |\xi|^2 \e^{2t})^{-\frac{\beta}{2}} \dd \xi \\
    + \sum_{\substack{\delta \ll 1 \\ \mathrm{dyadic}}} \delta^{-\beta} \sum_{1 \leq  c \ll \delta \e^t} \sum_{ 0 < |d| \ll |c|} \tfrac{\delta}{|c|} \e^{-t} \ll 1,
  \end{multline}
  where $\epsilon> 0$ is chosen as in Lemma \ref{lem:basicfundamentaldomain} to ensure $d \neq 0$. 

  For $d > 1$, we let fix $s > 0$ large enough so that, applying Lemma \ref{lem:crudecontraction},
  \begin{equation}
    \label{eq:sfixeddef}
    \int_{-1}^1 \alpha_d(gn(\xi)a(s))^\beta \dd \xi \leq \tfrac{1}{3} \alpha_d(g)^\beta + O( \int_{-1}^1 \alpha_{d-1}(gn(\xi)a(s))^\beta \dd \xi). 
  \end{equation}
  We write $t = ks + t_1$ with integer $k\geq 0$ and $0\leq t_1 < s$.
  By Lemma \ref{lem:rhotransbounds} we have
  \begin{equation}
    \label{eq:ttoksbound}
    \int_{-10}^{10} \alpha_d(gn(\xi) a(t))^\beta \dd \xi \ll \int_{-10}^{10} \alpha_d(gn(\xi) a(ks))^\beta \dd \xi.
  \end{equation}
  We now show that the right side of \eqref{eq:ttoksbound} is bounded by iterating the following.
  \begin{equation}
  \begin{split}
    \label{eq:contractioniteration}
&    \int_{-10}^{10} \alpha_d(gn(\xi) a(ks))^\beta \dd \xi \\
 &   \leq 2 \int_{-10}^{10} \int_{-1}^1 \alpha_d(gn(\xi) a((k-1)s)n(\xi_1)a(s))^\beta \dd \xi_1 \dd \xi \\
  &  \leq \tfrac{2}{3} \int_{-10}^{10} \alpha_d(gn(\xi) a((k-1)s))^\beta \dd \xi + O( \int_{-10}^{10} \alpha_{d-1}(gn(\xi) a(ks))^\beta \dd \xi) \\
  &  \leq \tfrac{2}{3} \int_{-10}^{10} \alpha_d(gn(\xi) a((k-1)s))^\beta \dd \xi + O( 1),
  \end{split}
  \end{equation}
  where the last bound follows by induction.
  The bound \eqref{eq:smallmoment} now follows from the case $k=0$, which is implied by Lemma \ref{lem:rhotransbounds}. 
  \end{proof}

\subsection{Proof of Theorem \ref{thm:escapeofmass}}
\label{sec:propproof}

Let $K>1$ and $\delta=\frac{c}{K}$ for some constant $c>0$ to be determined later. For this $\delta$ and sufficiently large $t=\log M$, we can find a finite sequence $\set{s_i}_{1\leq i\leq N}$ with $s_i=(1+\delta)s_{i+1}$ for $1\leq i\leq N-1$ and $T\leq s_N\leq 2T$ for some given $T>0$ sufficiently large. Then the sequence $\set{s_i}_{1\leq i\leq N}$ satisfies
        \begin{equation}
        \label{eq:sibound}
        \log M=s_1+\cdots+s_N, %\quad
        %s_N \textrm{ sufficiently large but bounded},
    \end{equation}
    \begin{equation}
      \label{eq:sumratio}
      \left(1-\frac{s_N}{s_i}\right)\del^{-1}s_i\leq s_{i+1}+\cdots+s_N\leq \del^{-1}s_i 
      \textrm{ for any }1\leq i\leq N-1.
  \end{equation}
    
We recall from Lemma \ref{lem:crudecontraction} that
    \eqlabel{eq:recalltrivialcontraction}{
        \int_{-1}^{1}\alpha_d\big(hn(\xi)a(s)\big)^{\beta}d\xi\leq  e^{(3\beta-2 + \epsilon ) s}\alpha_3(h)^{\beta}+e^{O(1)s}\alpha_{2}(h)^{\beta}}
    for any $\Gamma^3 h\in \Gamma^3\backslash G^3$ and sufficiently large $s$, and from Proposition \ref{prop:globalcontraction}
    \eqlabel{eq:recallnontrivialcontraction}{
        \int_{-1}^{1}\widetilde{\alpha}_d\big(hn(\xi)a(s)\big)d\xi\leq e^{-3(2-\beta - \epsilon) s}\widetilde{\alpha}(h)+e^{(K+O(1))s}\alpha_{2}(h)}
      for any $\Gamma h\notin\cE_{s,K}$ and sufficiently large $s$.
      Here and in what follows implied constants $\ll$ may depend on $K$ and $\delta$, but $O(\cdot)$ in the exponent are independent of $K$ and $\delta$.

      We note that from Lemma \ref{lem:alpha3supbound} and Proposition \ref{prop:smallmoment},
\begin{multline}
  \label{eq:boundedxi}
  \tfrac{1}{L} \int_{\e^{-(1 + \eta)t} < |\xi| \leq 10} \alpha_3(g_0 n(\xi)a(t))^\beta \dd \xi \\
  \ll \tfrac{1}{L} \e^{(\beta - \frac{2}{3} + \epsilon_1) (\frac{3\kappa}{\kappa + 1} + 3\eta)t} \int_{|\xi| \leq 10} \alpha_3(g_0 n(\xi)a(t))^{\frac{2}{3} - \epsilon_1} \dd \xi \ll_\epsilon \tfrac{1}{L} \e^{((3\beta - 2)(\frac{\kappa}{\kappa + 1} + \eta) + \epsilon)t}
\end{multline}
for any $\epsilon > 0$.
This contributes the second term on the right of (\ref{eq:escapeofmass}).

We have
\eqlabel{eq:Ztexpression}{L^{-1}\int_{10}^L \alpha_3\big(g_0n(\xi)a(t)\big)^{\beta}d\xi\leq 2\mathsf{Z}_t,}
where
$$\mathsf{Z}_t:=L^{-1}\int_{3}^{3L}\int \alpha_3\big(g_0n(\xi_0)n(\xi_1)a(s_1)\cdots n(\xi_N)a(s_N)\big)^{\beta}dm_I^{\otimes N}(\xi_1,\cdots,\xi_{N})d\xi_0$$
with $I = [-1,1]$ and $m_I$ is Lebesgue measure on $I$.
A similar bound holds for the integral over $-L \leq \xi \leq -10$, so we shall now prove that $\displaystyle\sup_{t>0}\mathsf{Z}_t<\infty$.
    
    For $1 \leq m \leq N-1$, let us define
    $$\overline{\Theta}_m:=\set{(\xi_0,\cdots,\xi_m)\in [3,3L]\times I^{m-1}: \Gamma^3 g_0n(\xi_0)n(\xi_1)a(s_1)\cdots n(\xi_m)a(s_m)\in \cE_{s_{m+1},K}},$$
    $$\Theta_m:=\overline{\Theta}_m\times I^{N-m}\subseteq [3,3L]\times I^{N-1},\ \Theta:=\bigcup_{m=1}^{N-1}\Theta_m.$$
    We have $\mathsf{Z}_t \leq \mathsf{Y}_t+\sum_{m=1}^{N-1}\mathsf{E}_{t,m}$, where
    $$\mathsf{Y}_t:=L^{-1}\int_{[3,3L]\times I^{N-1}\setminus\Theta} \alpha_3\big(g_0n(\xi_0)n(\xi_1)a(s_1)\cdots n(\xi_N)a(s_N)\big)^{\beta}dm_I^{\otimes N}(\xi_1,\cdots,\xi_{N}),$$
    $$\mathsf{E}_{t,m}:=L^{-1}\int_{\Theta_m} \alpha_3\big(g_0n(\xi_0)n(\xi_1)a(s_1)\cdots n(\xi_N)a(s_N)\big)^{\beta}dm_I^{\otimes N}(\xi_1,\cdots,\xi_{N})$$
    for $1\leq m\leq N-1$. 
    
    We shall first estimate $\mathsf{E}_{t,m}$. Write
    $$\mathsf{E}_{t,m}=L^{-1}\int_{\overline{\Theta}_m}J(\xi_0,\xi_1,\cdots,\xi_m)dm_I^{\otimes m}(\xi_1,\cdots,\xi_m),$$
    where
    $$J(\xi_0,\xi_1,\cdots,\xi_m):=\int_{I^{N-m}}\alpha_3\big(g_0n(\xi_1)a(s_1)\cdots n(\xi_N)a(s_N)\big)^{\beta}dm_I^{\otimes (N-m)}(\xi_{m+1},\cdots,\xi_{N})d\xi_0.$$
    
    Since $s_i\geq T$ for any $1\leq i\leq N$, we may apply \eqref{eq:recalltrivialcontraction} for $J(\xi_1,\cdots,\xi_m)$ repeatedly with $s=s_N,\cdots,s_{m+1}$, and get
    \eqlabel{eq:Badestimate1}{\begin{aligned}
        J(\xi_0,\xi_1,\cdots,\xi_m)&\leq e^{(3\beta-2 + \epsilon)(s_{m+1}+\cdots+s_N)}\alpha_{3}\big(g_0n(\xi_0)n(\xi_1)a(s_1)\cdots n(\xi_m)a(s_m)\big)^{\beta}\\&\quad+\sum_{i=m+1}^{N}e^{O(s_i)}e^{(3\beta-2)(s_{i+1}+\cdots+s_{N})}\alpha_{2}\big(g_0n(\xi_0)n(\xi_1)a(s_1)\cdots n(\xi_m)a(s_m)\big)^{\beta}\\&\leq (N-m)e^{(3\beta-2)(s_{m+2}+\cdots+s_{N})+O(s_{m+1})}\alpha_{3}\big(g_0n(\xi_0)n(\xi_1)a(s_1)\cdots n(\xi_m)a(s_m)\big)^{\beta}\\&\ll e^{\{(3\beta-2)\delta^{-1}+O(1)\}s_{m+1}}\alpha_{3}\big(g_0n(\xi_0)n(\xi_1)a(s_1)\cdots n(\xi_m)a(s_m)\big)^{\beta},
    \end{aligned}}
  using \eqref{eq:sumratio} in the penultimate inequality.
    
    Let $\varphi:=\alpha_3^{\beta}\cdot\mathds{1}_{\cE_{s_{m+1},K}}$. By Proposition \ref{prop:avoidanceMargulis} we have
    \eq{\begin{aligned}\label{eq:thetabarint}
        L^{-1}\int_{\widebar{\Theta}_m}\alpha_3&\big(g_0n(\xi_0)n(\xi_1)a(s_1)\cdots n(\xi_m)a(s_m)\big)^{\beta}dm_I^{\otimes m}(\xi_1,\cdots,\xi_m)d\xi_0\\&=L^{-1}\int_{[3,3L]\times I^{m-1}}\varphi\big(g_0n(\xi_0)n(\xi_1)a(s_1)\cdots n(\xi_m)a(s_m)\big)dm_I^{\otimes m}(\xi_1,\cdots,\xi_m)d\xi_0\\&\leq L^{-1}\int_{1}^{10L}\varphi(g_0n(\xi)a(s_1+\cdots+s_m))d\xi\ll e^{-(\mu K+O(1))s_{m+1}}.
    \end{aligned}}
    
    In combination with \eqref{eq:Badestimate1} it follows that
    \eqlabel{eq:Badestimate3}{\begin{aligned}\mathsf{E}_{t,m}&\ll e^{\{(3\beta-2)\delta^{-1}+O(1)\}s_{m+1}}  \\ & \quad \times L^{-1} \int_{\widebar{\Theta}_m}\alpha_3\big(g_0n(\xi_0)n(\xi_1)a(s_1)\cdots n(\xi_m)a(s_m)\big)^{\beta}dm_I^{\otimes m}(\xi_1,\cdots,\xi_m)\\&\ll e^{-\{\mu K-(3\beta-2)\delta^{-1}+O(1)\}s_{m+1}}
    \end{aligned}}
    for any $0\leq m\leq N-1$. Therefore, if 
    \eqlabel{eq:Condition1}{c\mu=\mu \delta K >3\beta -2}
    then we have
    \eqlabel{eq:Ebound}{\begin{aligned}
        \sum_{m=1}^{N-1}\mathsf{E}_{t,m} &\ll  \sum_{m=1}^{N-1}e^{-\{\mu K-(3\beta-2)\delta^{-1}+O(1)\}s_{m+1}}\\&< \sum_{m=1}^{\infty} e^{-(1+\del)^m\{\mu K-(3\beta-2)\delta^{-1}+O(1)\}T}\\&<\sum_{m=1}^{\infty} e^{-m\del\{\mu K-(3\beta-2)\delta^{-1}+O(1)\} T}\\&=\sum_{m=1}^{\infty} e^{-m\{(c\mu -(3\beta-2))+O(K^{-1})\} T}\ll 1
    \end{aligned}}
    for sufficiently large $K$.

    We now estimate $\mathsf{Y}_t$ using \eqref{eq:recallnontrivialcontraction}. Observe that for any $1\leq m\leq N$ and $(\xi_0, \xi_1,\cdots,\xi_{N})\in [3,3L]\times I^{N-1}\setminus\Theta$ we have 
    $$\Gamma^3 g_0n(\xi_0)n(\xi_1)a(s_1)\cdots n(\xi_m)a(s_m)\notin \cE_{s_{m+1},K}$$ from the construction of $\Theta$.
    We also have the following.
   \begin{lem}
    \label{lem:initialDiophantine}
    For $1 \leq |\xi| \leq 3L$, $\delta$ sufficiently small and $K$ sufficiently large, we have for $1 \leq \kappa < \tfrac{1}{2} K \delta$ that $g_0 n(\xi) \not\in \cE_{s_1, K}$.
  \end{lem}

  \begin{proof}
    Let $\gamma = \gamma(g_0 n(\xi)n(\xi_1)a(s_1))= \left(
    \begin{pmatrix}
      * & * \\
      c_j & d_j 
    \end{pmatrix}
    \right)_{\!\! j=1,2,3}$.
    We may assume that $c_j \neq 0$, then $|\xi + \xi_1 + \frac{d_j}{c_jx_j}| \ll |c_j|^{-1} \e^{-s_1}$, so in particular $\xi + \xi_1 \geq 2$ implies that $d_j \neq 0$.
    Now by the Diophantine condition on $\tfrac{x_i}{x_j}$, for $c_ic_j \neq 0$, we have
    \begin{equation}
      \label{eq:initialexceptional}
      L^{-\kappa} \e^{-(2\kappa + 1)s_1} \ll |\tfrac{d_i}{c_i x_i } - \tfrac{d_j}{c_j x_j}|.
    \end{equation}
    Since $t \leq \delta^{-1} s_1$, this implies $|\tfrac{d_i}{c_i x_i } - \tfrac{d_j}{c_j x_j}| > \e^{-\frac{1}{2} K s_1}$ as long as $\kappa <  \tfrac{1}{2} K \delta$ and $K$ and $\delta$ are sufficiently large and small respectively.
  \end{proof}
    
    
    Applying \eqref{eq:recallnontrivialcontraction} with $s=s_N,\cdots,s_1$ repeatedly, we get
    \eqlabel{eq:Goodestimate1}{\begin{aligned}
        \mathsf{Y}_t&<L^{-1}\int_{[3,3L]\times I^{N-1}\setminus\Theta}\alpha_3\big(g_0n(\xi_0)n(\xi_1)a(s_1)\cdots n(\xi_N)a(s_N)\big)^{\beta}dm_I^{\otimes N}(\xi_1,\cdots,\xi_{N})d\xi_0\\&\leq L^{-1}\int_{[3,3L]\times I^{N-1}\setminus\Theta} \widetilde{\alpha}\big(g_0n(\xi_0)n(\xi_1)a(s_1)\cdots n(\xi_N)a(s_N)\big)dm_I^{\otimes N}(\xi_1,\cdots,\xi_{N})d\xi_0\\&\leq e^{-3(2-\beta - \epsilon) t}L^{-1}\int_{3}^{3L}\widetilde{\alpha}(g_0n(\xi_0))d\xi_0%\\&\quad+e^{(\frac{d-1}{l-1}K_0+O(1))s_1-d(2-\beta)(s_{2}+\cdots+s_{N})}L^{-1}\int_{-L}^{L}\int_{I}\max_{1\leq k\leq d}\alpha_{d-1}\big(\pi_{d-1,k}(g)n(\xi_1)a(s_1)\big)^{\beta}d\xi_1
        \\&\qquad+\sum_{m=1}^{N}e^{(K+O(1))s_m-3(2-\beta)(s_{m+1}+\cdots+s_{N})}\\&\qquad\qquad\times L^{-1}\int_{3}^{3L}\int_{I^{m}}\alpha_{2}\big(g_0n(\xi_0)n(\xi_1)a(s_1)\cdots n(\xi_m)a(s_m)\big)dm_I^{\otimes m}(\xi_1,\cdots,\xi_{m})d\xi_0\\&\ll e^{-3(2-\beta-\epsilon) t}L^{-1}\int_{3}^{3L}\widetilde{\alpha}(g_0n(\xi_0))d\xi_0 +\sum_{m=1}^{N}e^{(K+O(1))s_m-3(2-\beta)(s_{m+1}+\cdots+s_{N})}.
    \end{aligned}}
    We choose a sufficiently small $\epsilon>0$ and let $N'$ be the largest integer such that $s_{N'}\geq \epsilon^{-1}s_N$. Note that $N-N'< \delta^{-1}\log\frac{1}{\eps}$, since $s_{N'}=(1+\delta)^{N-N'}s_N$. If $1\leq m\leq N'$, then $$s_{m+1}+\cdots+s_{N}\geq \left(1-\frac{s_N}{s_{N'}}\right)\delta^{-1}s_m\geq(1-\epsilon)\delta^{-1}s_m.$$
    Hence, if 
    \eqlabel{eq:Condition2}{c=K\delta<3(1-\epsilon)(2-\beta)} then we have
    \eqlabel{eq:exponentialsum0}{\begin{aligned}
        \sum_{m=1}^{N'}e^{(K+O(1))s_m-3(2-\beta)(s_{m+1}+\cdots+s_{N})}&\leq \sum_{m=1}^{N'}e^{\{(K+O(1))-3(1-\epsilon)(2-\beta)\delta^{-1}\}s_m}\\&\ll\sum_{m=1}^{\infty}e^{\{(K+O(1))-3(1-\epsilon)(2-\beta)\delta^{-1}\}(1+\delta)^mT}\\&\leq \sum_{m=1}^{\infty}e^{\{(c-3(1-\epsilon)(2-\beta))+O(K^{-1)}\}T}<\infty
    \end{aligned}}
    uniformly for sufficiently large $K$.  
     
    Now we fix sufficiently large $K$ so that \eqref{eq:Ebound} and \eqref{eq:exponentialsum0} hold. Since $\mu > \tfrac{1}{3}$, there exists $\beta>1$ such that $\mu^{-1}(3\beta-2)<3(2-\beta)$. Thus for sufficiently small $\epsilon>0$ we can also find a constant $c$ with $$\mu^{-1}(3\beta-2)<c<(1-\epsilon)3(2-\beta)$$ and $\delta=\frac{c}{K}$ satisfying both \eqref{eq:Condition1} and \eqref{eq:Condition2}. Note that this condition also ensures the assumption $\kappa<\frac{1}{2}K\delta=\tfrac{c}{2}$ in Lemma \ref{lem:initialDiophantine}.
    
    It is also easy to see
    \eqlabel{eq:exponentialsum2}{\begin{aligned}
        \displaystyle\sum_{m=N'+1}^{N}e^{(K+O(1))s_m-3(2-\beta)(s_{m+1}+\cdots+s_{N})}&\leq (N-N')e^{(2K+O(1))s_N}\\& <\delta^{-1}\log\frac{1}{\epsilon}e^{(2K+O(1))s_N}<\infty.
    \end{aligned}}
  
    Combining \eqref{eq:Goodestimate1}, \eqref{eq:exponentialsum0}, and \eqref{eq:exponentialsum2} with Lemma \ref{lem:initialtildealphaaverage} below, $\mathsf{Y}_t$ is uniformly bounded for any large $t$. In combination with \eqref{eq:Ebound}, it follows that
    $$\mathsf{Z}_t=\mathsf{Y}_t+\sum_{m=1}^{N}\mathsf{E}_{t,m}$$
    is also uniformly bounded for any large $t$, hence $\displaystyle\sup_{t>0}\mathsf{Z}_t<\infty$.
    This completes the proof of Theorem \ref{thm:escapeofmass} once we have the following.

  \begin{lem}
    \label{lem:initialtildealphaaverage}
    We have
    \begin{equation}
      \label{eq:initialtildealphaaverage}
      \tfrac{1}{L} \int_3^{3L} \tilde{\alpha}(g_0 n(\xi)) \dd \xi \ll L^{2\kappa -1}. 
    \end{equation}
  \end{lem}

  \begin{proof}
    Let $m_j$ be the integer closest to $x_j \xi$.
    Then either $\gamma_j(g_0n(\xi)) =
    \begin{pmatrix}
      1 & -m_j \\
      0 & 1
    \end{pmatrix}
    $ or $\gamma_j(g_0n(\xi)) =
    \begin{pmatrix}
      0 & -1 \\
      1 & -m_j
    \end{pmatrix}
    $.
    %depending on whether $x_j \sqrt{(\xi - m_j)^2 + 1} \geq 1$ or $x_j \sqrt{(\xi - m_j)^2 + 1} < 1$.
    We may assume that for some $i\neq j$ we have the second case
    since otherwise $\Delta_3(g_0n(\xi))^{-1} =1$ and $\alpha_3(g_0n(\xi)) \asymp 1$.

    Now $\Delta_3(g_0 n(\xi))^{-1} \ll \max_{i\neq j} |\tfrac{m_i}{x_i} - \tfrac{m_j}{x_j}|^{-2}$, so we may bound the left side of \eqref{eq:initialtildealphaaverage} by
    \begin{equation}
      \label{eq:mimjsum}
      \ll \tfrac{1}{L} \sum_{1\leq i\neq j\leq 3} \sum_{3\leq m_i,m_j \leq 3L} |\tfrac{m_i}{x_i} - \tfrac{m_j}{x_j}|^{-2} . 
    \end{equation}
    We let $\tfrac{r}{q}$ be the best rational approximation to $\frac{x_i}{x_j}$ with denominator $\leq 3L$, so that $|\tfrac{x_i}{x_j} - \tfrac{r}{q}| \leq \tfrac{1}{3Lq}$ and $|\tfrac{m_i}{m_j} - \tfrac{x_i}{x_j}| \gg |\tfrac{m_i}{m_j} - \tfrac{r}{q}|$ for $\tfrac{m_i}{m_j} \neq \tfrac{r}{q}$.

    Set $h = m_iq - m_j r$.
    We note that $m_j$ is determined modulo $q$ from $h$, and then $m_i$ is determined from $h$ and $m_j$.
    The contribution from $h \neq 0$ is given by
    \begin{equation}
      \label{eq:mimjhnot0}
      \tfrac{1}{L} q^2 \sum_{h\neq 0} h^{-2} \sum_{\substack{ m \leq 3L \\ m \equiv h \bar r \bmod q}} 1 \ll q. 
    \end{equation}
    For $m_i q - m_j r = 0$, ie $m_j = mq$, $m_i = mr$, $m \ll \tfrac{L}{q}$, we note that since $\tfrac{x_i}{x_j}$ is $\kappa$-Diophantine, $|\tfrac{m_i}{x_i} - \tfrac{m_j}{x_j}| \geq (mq)^{-\kappa}$.
    The contribution to \eqref{eq:initialtildealphaaverage} is therefore $\tfrac{1}{L} q^{2\kappa}$. 
  \end{proof}

\section{Avoidance estimate}
\label{sec:avoidance}

In this section we prove the following avoidance estimate, which is used in the proof of Theorem \ref{thm:escapeofmass}, see \eqref{eq:thetabarint}.
\begin{prop}\label{prop:avoidanceMargulis}
  There exists $\kappa_0 > 1$ and $\mu > \tfrac{1}{3}$ such that for $1 \leq \kappa < \kappa_0$ there is $\beta = \beta(\kappa) > 1$ with the following property.
  For $g_0 = \left(\begin{pmatrix}    x_j^{\frac{1}{2}} & 0 \\
    0 & x_j^{-\frac{1}{2}}
  \end{pmatrix}
  \right)_{\!\! j=1,2,3}$ with $\tfrac{x_i}{x_j}$ $\kappa$-Diophantine, we have
  \begin{equation}
    \label{eq:avoidanceMargulis}
    \int_{3 < |\xi| \leq  L} \alpha_3\big(gn(\xi)a(t)\big)^{\beta}\mathds{1}_{\cE_{s,K}}(\Gamma gn(\xi)a(t))d\xi\ll e^{-(\mu K+O(1))s}L
  \end{equation}
  uniformly for $K, s \geq 0$. 
\end{prop}

\subsection{Avoidance in $\Gamma^2 \backslash G^2$}
\label{sec:rank2avoidance}

Towards proving Proposition \ref{prop:avoidanceMargulis}, we prove the following in $\Gamma^2 \backslash G^2$.
For $g =\left(
\begin{pmatrix}
  * & * \\
  c_j & d_j 
\end{pmatrix}
\right)_{\!\! j=1,2}
\in G^2$, we let $\Delta_2(g) = \min(1, |\tfrac{d_1}{c_1} - \tfrac{d_2}{c_2}|)$ and we set
\begin{multline}
  \label{eq:cE'def}
  \cE'_{s,K} = \{ \Gamma^2 g \in \Gamma^2 \backslash G^2: \Delta_2(\gamma g) \leq \e^{-Ks}\textrm{ for some } \gamma \in \Gamma^2\\
  \textrm{ such that } \max_{j=1,2} \rho(\gamma_j g_j) < 10 \e^s \}. 
\end{multline}
We note the similarty between $\cE'_{s,K}$ and $\cE_{s,K}$ in (\ref{eq:tildecEdef}) part from the lack of the lower bound for $\rho$ in (\ref{eq:cE'def}).
Roughly speaking, $\cE'_{s,K}$ is obtained from $\cE_{s,K}$ by removing the coordinate $j$ such that $\rho(\gamma_j g_j) > \tfrac{1}{10} \e^{-s}$. 
\begin{prop}
  \label{prop:rank2avoidance}
  Let $g_0 = \left(\begin{pmatrix}    x_j^{\frac{1}{2}} & 0 \\
    0 & x_j^{-\frac{1}{2}}
  \end{pmatrix}
  \right)_{\!\! j=1,2}$ with $\tfrac{x_1}{x_2}$ $\kappa$-Diophantine.
  Then for any $\beta < 1 + \kappa^{-2}$ and $1\leq L \ll \e^t$, we have
  \begin{equation}
    \label{eq:rank2moment1}
    \int_{1 \leq |\xi| \leq L} \alpha_2(g_0 n(\xi) a(t))^\beta \mathbbm{1}_{\cE'_{s,K}}(gn(\xi)a(t)) \dd \xi \ll L \e^{-(\kappa^{-2} K + O(1))s} .
  \end{equation}
\end{prop}

The proof of Proposition \ref{prop:rank2avoidance} requires the following lemmas.

For $1\leq C_j \ll \e^t$, define
\begin{multline}
  \label{eq:cCdef2}
  \cC_{t,\eta, L}(C_1, C_2) = \{(c_j, d_j)_{j=1,2} : \gcd(c_j,d_j) = 1,\ C_j \leq |c_j| < 2 C_j,\\
  0 < |d_j| \leq 2x_j LC_j,\ |\tfrac{d_1}{c_1x_1} - \tfrac{d_2}{c_2x_2}| \leq \eta \e^{-2t}\}. 
\end{multline}

\begin{lem}
  \label{lem:cCcount}
  Suppose that $\tfrac{x_1}{x_2}$ is $\kappa$-Diophantine.
  \begin{enumerate}[{\rm (i)}]
  \item We have $|\cC_{t, \eta, L}(C_1, C_2)| = 0$ unless $(L C_1C_2)^{-\kappa} \ll \eta C_1C_2 \e^{-2t}$. 
    
  \item
    We have
    \begin{equation}
      \label{eq:cCcount}
      |\cC_{t, \eta, L}(C_1,C_2)|\ll_\epsilon
      %(\eta C_1C_2\e^{-2t} +(LC_1C_2)^{-\frac{1}{\kappa}})(LC_1C_2)^{1+\epsilon}
      \begin{cases}
        (\eta C_1C_2\e^{-2t} +(LC_1C_2)^{-\frac{1}{\kappa}})(LC_1C_2)^{1+\epsilon} \\
        (\eta C_1C_2\e^{-2t} +(\eta C_1C_2\e^{-2t})^{\kappa^{-2}})(LC_1C_2)^{1+\epsilon}
        % \eta^{\kappa^{-2}} (C_1C_2)^{1+ \kappa^{-2} + \epsilon} L^{1 + \epsilon} \e^{-2\kappa^{-2}t}.
      \end{cases}
    \end{equation}
    for any $\epsilon > 0$.
  \end{enumerate}
\end{lem}

\begin{proof}
  We first note that since $c_1d_2 \neq 0$,
  \begin{equation}
    \label{eq:Diophantinelowerbound}
    |c_1d_2 \tfrac{x_1}{x_2} - c_2d_1| \geq |c_1d_2|^{-\kappa} \gg (LC_1C_2)^{-\kappa}.
  \end{equation}
  and (i) follows.
  
  For (ii), let $c = c_1d_2$ and $d = c_2d_1$.
  Let $1\leq q \leq LC_1C_2$ be so that $|\tfrac{x_2}{x_1} - \tfrac{r}{q}| \leq \tfrac{1}{LC_1C_2 q}$ for some integer $r$ coprime to $q$.
  Since $\frac{x_2}{x_1}$ is $\kappa$-Diophantine, we have $q^{\kappa}\gg LC_1C_2$.
  
  From $|c\tfrac{x_1}{x_2} - d| \ll \eta C_1 C_2 \e^{-2t}$, we have
  \begin{equation}
    \label{eq:hdef}
    h := cr - dq \ll \eta q C_1C_2 \e^{-2t}  + 1. 
  \end{equation}
    By assertion (i), we may assume $$\eta C_1C_2 \e^{-2t} \gg (LC_1C_2)^{-\kappa}\gg q^{-\kappa^2},$$
    since otherwise $\cC_{t, \eta, L}(C_1, C_2)$ is empty. 
  
  We note that $c \equiv h \overline{r} \bmod q$, and $d = \tfrac{cr - h}{q}$, so
  \begin{equation}
  \begin{split}
    \label{eq:cCcount1}
    |\cC_{t, \eta, L}(C_1, C_2)| & \ll \sum_{|h| \ll \eta q C_1 C_2 \e^{-2t} +1} \sum_{\substack{|c| \leq LC_1C_2 \\ c \equiv h\overline{r} \bmod q}} \tau(c) \tau(\tfrac{cr - h}{q})\\ & \ll \tfrac{1}{q} (\eta qC_1C_2 \e^{-2t} +1)(LC_1C_2)^{1+\epsilon},
  \end{split}
  \end{equation}
  where $\tau$ is the divisor function.
\end{proof}

Combined with Lemma \ref{lem:cCcount}, \cite[Theorem~1]{LMW2023a} quickly gives the following.
\begin{lem}
  \label{lem:rank2moments1}
  Let $g_0 = \left(\begin{pmatrix}    x_j^{\frac{1}{2}} & 0 \\
    0 & x_j^{-\frac{1}{2}}
  \end{pmatrix}
  \right)_{\!\! j=1,2}$ with $\tfrac{x_1}{x_2}$ $\kappa$-Diophantine.
  Then for any $1\leq L \ll \e^t$, $\beta < 1 + \kappa^{-2}$ and $\epsilon > 0$, we have
  \begin{equation}
    \label{eq:rank2moment1B}
    \int_{1\leq |\xi| \leq L} \alpha_2(g_0 n(\xi) a(t))^\beta \dd \xi \ll L.
  \end{equation}
\end{lem}

\begin{proof}
  We use Lemma \ref{lem:cCcount} to bound the contribution of $\xi$ such that $\alpha_2(g_0n(\xi)a(t)) > \e^{\epsilon_1 t}$.
  Since the indicator function of the region where $\alpha_2(g) \leq \e^{\epsilon_1 t}$ has a smooth compactly supported majorant with Sobolev norm $\ll \e^{O(\epsilon_1 t)}$, the remaining portion is bounded by means of Proposition \ref{prop:effectiveequidistribution}, noting that for $\beta < 2$,
  \begin{equation}
    \label{eq:alpha2average}
    \int_{\Gamma^2 \backslash G^2} \alpha_2(g_0 n(\xi) a(t))^\beta \dd \xi \ll 1
  \end{equation}
  and that by taking $\epsilon_1$ sufficiently small, the remainder term on the right of \eqref{eq:effectiveequidistribution} is also bounded. 

  We dyadically decompose according to $\tfrac{1}{2} \delta_j^{-1} < \alpha_1(g_j n(\xi) a(t))\leq \delta_j^{-1}$. 
  We note that from Lemma \ref{lem:basicfundamentaldomain}, $\gamma(g_0n(\xi)a(t)) = \left(\begin{pmatrix}    * & * \\
    c_j & d_j
  \end{pmatrix}
  \right)_{\!\! j=1,2}$ satisfies $0 < |c_j| \ll \delta_j \e^t$, $d_j \neq 0$, and $|\xi + \tfrac{d_j}{c_jx_j}| \ll \tfrac{\delta_j}{|c_j|\e^t}$, so in particular $|\tfrac{d_1}{c_1 x_1} - \tfrac{d_2}{c_2 x_2}| \ll (\tfrac{\delta_1}{|c_1|} + \tfrac{\delta_2}{|c_2|}) \e^{-t}.$
  Moreover, given such a $\gamma$, the measure of the associated $\xi$ is $O ( \min ( \tfrac{\delta_j}{c_j\e^t}))$ and therefore we obtain the bound
  \begin{multline}
    \e^{-t}\sum_{\substack{\delta_1,\delta_2 \ll \e^{-\epsilon_1 t} \\ \delta_j \ll 1\\ \mathrm{dyadic}}} (\delta_1\delta_2)^{-\beta} \sum_{\substack{1\leq C_j \ll \delta_j \e^t \\ \mathrm{dyadic}}} \min( \tfrac{\delta_j}{C_j}) |\cC_{t,(\delta_1 C_2 + \delta_2 C_1)\e^t,L}(C_1, C_2)|.
  \end{multline}

  We bound $\cC$ using Lemma \ref{lem:cCcount}.
  For $\epsilon_2 > 0$, the first term on the right of \eqref{eq:cCcount} contributes
  \begin{equation}
    \label{eq:cCcountfirstterm}
    L^{1+\epsilon_2} \e^{-2t} \sum_{\substack{\delta_1,\delta_2 \ll \e^{-\epsilon_1 t}\\ \delta_j \ll 1 \\ \mathrm{dyadic}}} (\delta_1\delta_2)^{1 -\beta} \sum_{\substack{1\leq C_j \ll \delta_j \e^t \\ \mathrm{dyadic}}} (C_1C_2)^{2 + \epsilon_2} \\
    \ll L^{1 + \epsilon_2} \e^{-(2 -\beta - 2\epsilon_2) \epsilon_1 t}.
  \end{equation}
  Using $L \ll \e^t$, $\beta < 2$ and taking $\epsilon_2$ sufficiently small makes this $O(L)$.
  
  Now using Lemma \ref{lem:cCcount} (i), the second term on the right of \eqref{eq:cCcount} contributes 
  \begin{multline}
    \label{eq:cCcountsecondterm}
    L^{1 + \epsilon_2 } \e^{-(1 + \kappa^{-2})t} \sum_{\substack{ \delta_1\delta_2 \ll \e^{-\epsilon_1 t} \\ \delta_j \ll 1\\ \mathrm{dyadic}}} (\delta_1 \delta_2)^{- \beta}
  \sum_{\substack{C_j \ll \delta_j \e^t \\\mathrm{dyadic}}} \min(\frac{\delta_j}{C_j}) \max( \frac{\delta_j}{C_j})^{\kappa^{-2}}(C_1C_2)^{1 + \kappa^{-2} + \epsilon_2} \\
    \ll L^{1 + \epsilon_2}\e^{2\epsilon_2 t} \sum_{\substack{ \delta_1\delta_2 \ll \e^{-\epsilon_1 t} \\ \delta_j \ll 1\\ \mathrm{dyadic}}} (\delta_1 \delta_2)^{1 + \kappa^{-2}- \beta}.
  \end{multline}
  Since $\beta < 1 + \tfrac{1}{\kappa^2}$, this is $O(L)$ by making $\epsilon_2$ sufficiently small.  
\end{proof}

\begin{proof}[Proof of Proposition \ref{prop:rank2avoidance}]
When $s$ is small, we apply Proposition \ref{prop:effectiveequidistribution}.
By Lemma \ref{lem:rank2moments1} we may restrict to the set $\cS$ of $\alpha_2(g) \leq \e^{\epsilon_1t}$.
Now for $g = \left(
\begin{pmatrix}
  * & * \\
  0 & * 
\end{pmatrix}
\begin{pmatrix}
  \cos \theta_j & -\sin \theta_j \\
  \sin \theta_j & \cos \theta_j
\end{pmatrix}
\right)_{\!\! j=1,2}$, $\kappa_2(g)\leq \e^{-Ks}$ implies $|\theta_1 - \theta_2| \ll \e^{-Ks}$.
Moreover, for any $g = (g_1,g_2) \in G^2$, the number of $\gamma = (\gamma_1, \gamma_2) \in \Gamma^2$ such that $\rho(\gamma_jg_j) \leq 10 \e^s$ is $\ll \e^{O(s)}$.
Therefore, we may find smooth majorants $\varphi_1$ and $\varphi_2$ for the indicator functions of $\cS$ and $\cE'_{s,K}$ so that their Sobolev norms are $\ll \e^{O(\epsilon_1 t)}$ and $\ll \e^{O(K s)}$ respectively and
\begin{equation}
  \label{eq:varphijint}
  \int_{\Gamma^2\backslash G^2} \varphi_1(g)\varphi_2(g) \dd \mu(g) \ll \e^{-(K + O(1))s}.
\end{equation}
If $st^{-t}$ is sufficiently small, the remainder term from Proposition \ref{prop:effectiveequidistribution} is also $\ll\e^{-(K + O(1))s})$.

For the complementary range of $s$, we dyadically decompose as in the proof of Lemma \ref{lem:rank2moments1}, obtaining the bound
\begin{equation}
  \label{eq:avoidancedyadicdecomp}
  \ll \e^{-t}\sum_{\substack{\delta_j \ll \e^s\\ \mathrm{dyadic}}}(\delta_1 \delta_2)^{-\beta} \sum_{\substack{1\leq C_j \ll \delta_j \e^t \\ \mathrm{dyadic}}} \min(\tfrac{\delta_j}{C_j}) |\cC_{t,\e^{-Ks},L}(C_1,C_2)|.
\end{equation}
By Lemma \ref{lem:cCcount} this is bounded by
\begin{equation}
  \label{eq:avoidancedyadicdecomp1}
  \ll \e^{-\kappa^{-2}Ks} L^{1 + \epsilon_2} \e^{2\epsilon_2 t} \sum_{\substack{\delta_j \ll \e^s\\ \mathrm{dyadic}}}(\delta_1 \delta_2)^{1 + \kappa^{-2}-\beta},
\end{equation}
so the required bound follows by taking $\epsilon_2$ sufficiently small. 
\end{proof}
 

\subsection{Avoidance in $\Gamma^3 \backslash G^3$}
\label{sec:avoidanceproof}


We first decompose the exceptional set $\cE_{s,K}$ as follows:
\begin{equation}
  \label{eq:cEdecomp}
  \cE_{s,K}\subseteq\bigcup_{\substack{1\leq i\neq j \leq 3 \\ 1\leq k \leq 3}}\cE_{s,K}(i,j, k),
\end{equation}
where
\begin{multline}
  \label{eq:decompparts}
  \cE_{s,K}(i, j):=\{\Gamma^3 g\in \Gamma^3 \backslash G^3: 
  \Delta_{i,j}(\gamma g)^2 \leq\Delta_{j}(\gamma g)< \e^{-Ks} \\
  \textrm{ for some } \gamma \in \Gamma^3 \textrm{ such that } \tfrac{1}{10} \e^{-s}\leq \rho(\gamma_k g_k),\ \max_{j=1,2,3} \rho(\gamma_j g_j) \leq \tfrac{1}{10} \e^{s}.\}
  \end{multline}
with
\begin{equation}
  \label{eq:Deltadef2}
  \Delta_{j}\left(\left(\begin{pmatrix}    a_k & b_k\\
    c_k & d_k
  \end{pmatrix}
  \right)_{k=1,2,3}\right)= \prod_{i \neq j}|\tfrac{d_{i}}{c_{i}}-\tfrac{d_{j}}{c_{j}}|
\end{equation}
and
\begin{equation}
  \label{eq:Deltaijdef}
  \Delta_{ij}\left(\left(\begin{pmatrix}    a_k & b_k\\
    c_k & d_k
  \end{pmatrix}
  \right)_{k=1,2,3}\right) = |\tfrac{d_i}{c_i} - \tfrac{d_j}{c_j}|. 
\end{equation}
Without loss of generality, we may assume $i=2$ and $j=3$ and estimate the contribution of $\cE_{s,K}(2,3,k)$.

When $k = 1$, we have that
\begin{multline}
  \label{eq:rank2reduction1}
  \int_{3\leq |\xi| \leq L} \alpha_3(g_0n(\xi) a(t))^\beta \mathbbm{1}_{\cE_{s,K}(2,3,1)}(g_0n(\xi)a(t)) \dd \xi \\
  \ll \e^{O(s)} \int_{1\leq |\xi| \leq L} \alpha_2(\pi_1(g_0)n(\xi) a(t))^\beta \mathbbm{1}_{\cE'_{s,K/2}}(\pi_1(g_0)n(\xi)a(t)) \dd \xi,
\end{multline}
where $\pi_1: \Gamma^3 \backslash G^3 \to \Gamma^2 \backslash G^2$ is given by $\pi_1(g_1,g_2, g_3)= (g_2, g_3)$.
Now \eqref{eq:avoidanceMargulis} with $\cE_{s,K}$ replaced with $\cE_{s,K}(2,3,1)$ follows from Proposition \ref{prop:rank2avoidance} with $\mu < \tfrac{1}{2\kappa^2}$. 

When $k = 2$ (the case $k=3$ is similar), we have
$$\cE_{s,K}(2,3,2) \subset \bigcup_{\substack{K_1 + K_2 \geq K-1 \\ K_2 \geq K_1 -1}} \cE_{s,K_1,K_2}(2,3,2),$$
where $K_1, K_2$ are integers and
\begin{multline}
  \label{eq:etadef}
  \cE_{s,K_1,K_2}(2,3,2) = \{ g \in \Gamma^3 \backslash G^3 : \e^{-(K_i +1)s} < \Delta_{i3}(\gamma g) \leq \e^{-K_i s} \\
  \textrm{ for some } \gamma \in \Gamma^3 \textrm{ such that } \tfrac{1}{10} \e^{-s}\leq  \rho(\gamma_2 g_2),\ \max_{j=1,2,3} \rho(\gamma_j g_j) \leq \tfrac{1}{10} \e^{s}\}.
\end{multline}
Let us consider $K_2 \leq AK$ for some $A > 0$.
Let $p < \tfrac{1 + \kappa^{-2}}{\beta} < 2$.
By Holder's inequality, we have
\begin{multline}
  \label{eq:holder1}
  \sum_{\substack{K_1 + K_2 \geq K-1 \\ K_2 \leq AK}} \int_{3 \leq |\xi| \leq L} \alpha_2(\pi_2(g_0)n(\xi) a(t))^\beta \chi_{K_1}(g_0n(\xi)a(t)) \chi_{K_2}(g_0n(\xi)a(t)) \dd \xi \\
  \leq \sum_{\substack{K_1 + K_2 \geq K-1 \\ K_2 \leq AK}}\left( \int_{3\leq |\xi| \leq L} \alpha_2(\pi_2(g_0)n(\xi) a(t))^{p \beta} \tilde{\chi}_{K_1} (\pi_2(g_0)n(\xi) a(t))\dd \xi\right)^{\tfrac{1}{p}} \\
  \times \left(\int_{3 \leq |\xi| \leq L} \tilde{\chi}_{K_2}(\pi_1(g_0)n(\xi) a(t))\right)^{1 - \frac{1}{p}},
\end{multline}
with $\chi_{K_i}$ the indicator functions of the sets
\begin{multline}
  \label{eq:etasets}
  \{\Gamma^3 g\in \Gamma^3 \backslash G^3: \e^{-(K_i + 1)s} < \Delta_{i,3}(\gamma g) \leq \e^{-K_i s} \\
  \textrm{ for some } \gamma \in \Gamma^3 \textrm{ such that }
  \max_{1\leq k\leq 3}\rho(\gamma_k g_k)\leq 10 \e^s\}
\end{multline}
and $\tilde \chi_{K_i}$ is the indicator function of the associated projections of the sets \eqref{eq:etasets} to $\Gamma^2 \backslash G^2$ under $(g_1,g_2,g_3) \mapsto (g_i,g_3)$.
By Proposition \ref{prop:rank2avoidance}, \eqref{eq:holder1} is bounded by
\begin{equation}
  \label{eq:dyadicetasum}
  L \e^{O(s)}\sum_{\substack{K_1 + K_2 \geq K-1 \\ K_2 \leq AK}} \e^{-\frac{1}{\kappa^2 p}K_1 s} \e^{-\frac{1}{\kappa^2}(1 - \tfrac{1}{p})K_2 s }
  \ll L\e^{-((\frac{1}{\kappa^2 p} - \frac{A}{\kappa^2} (\frac{2}{p} - 1) )K + O(1))s}.
\end{equation}

Now consider the case $K_2 \geq AK$.
We let $\hat\pi_j : \Gamma^3 \backslash G^3 \to \Gamma\backslash G $ be given by $\hat\pi_j(g_1,g_2,g_3) = g_j$.
We let $\delta_1 > 0$ be so that $2\beta - \delta_1 < 1 + \tfrac{1}{\kappa^2}$, and we let $\tfrac{1}{p_1} + \tfrac{1}{q_1} = 1$ be so that $\beta_1 = 2p_1(\beta - \delta_1) < 1 + \tfrac{1}{\kappa^2}$ and $\beta_2 = q_1 \delta_1 < 1 + \tfrac{1}{\kappa^2}$.
Applying Holder we obtain the bound
\begin{multline}
  \label{eq:verysmalleta2}
  \sum_{\substack{K_1 + K_2 \geq K-1 \\ K_2 \geq AK}} \left(\int_{1\leq |\xi| \leq L} \alpha_1(\hat\pi_1(g_0)n(\xi) a(t))^{\beta_1}\dd \xi\right)^{\frac{1}{2p_1}} \\
    \times \left(\int_{1\leq |\xi| \leq L} \alpha_1(\hat\pi_3(g_0)n(\xi) a(t))^{\beta_1} \chi_{K_2}(g_0n(\xi)a(t)) \dd \xi\right)^{\frac{1}{2p_1}} \\
  \times \left(\int_{1\leq |\xi| \leq L} \alpha_2(\pi_2(g_0)n(\xi) a(t))^{\beta_2} \chi_{K_1}(g_0n(\xi)a(t)) \dd \xi\right)^{1 - \frac{1}{p_1}},
\end{multline}
and applying Proposition \ref{prop:rank2avoidance}, we obtain
\begin{equation}
  \label{eq:verysmalleta2eta1sum}
  L \e^{O(s)}\sum_{\substack{K_1 + K_2 \geq K-1 \\ K_2 \geq AK}}  \e^{-\frac{1}{2\kappa^2p_1}K_2s} \e^{-\frac{1}{\kappa^2q_1}K_1s} \\
  \ll L \e^{-(\frac{1}{\kappa^2}(1 - \frac{1}{p_1}) + \frac{A}{\kappa^2}(\frac{3}{2 p_1} - 1)K + O(1))s}
\end{equation}
assuming $p_1 < \tfrac{3}{2}$. 
We note that if one takes $\kappa$ and $\beta$ arbitarily close to $1$, then $q_1$ can be taken arbitarily large.
Then taking $A$ sufficiently large, one can then take $p$ sufficiently close to $2$ so that \eqref{eq:avoidanceMargulis} holds with $\mu > \tfrac{1}{3}$.
% optimize for kappa_0??

\section*{Data access statement}
No new data were generated or analysed during this study.

\bibliographystyle{amsalpha}
\bibliography{references}
\def\cprime{$'$} \def\cprime{$'$} \def\cprime{$'$}
\providecommand{\bysame}{\leavevmode\hbox to3em{\hrulefill}\thinspace}
\providecommand{\MR}{\relax\ifhmode\unskip\space\fi MR }
% \MRhref is called by the amsart/book/proc definition of \MR.
\providecommand{\MRhref}[2]{%
  \href{http://www.ams.org/mathscinet-getitem?mr=#1}{#2}
}

\end{document}
